\documentclass[11pt,a4paper]{article}
\begin{document}

% 제목 블록
\begin{center}
  \vspace*{2cm}
  {\Huge \textbf{중앙은행과 통화 정책: 금융 위기}} \\
  \vspace{0.5cm}
  {\Large \textbf{모듈 4: 금융 위기 (파트 1/2)}} \\
  \vspace{0.3cm}
  {\large Professor Ralf Meisenzahl} \\
  \vspace{0.2cm}
  {\normalsize FIN-571: Financial Institutions \& Monetary Policy} \\
  \vspace{1cm}
  \today
\end{center}

\newpage

\section*{개요}
이 문서는 '중앙은행과 통화 정책' 모듈 4의 첫 번째 부분으로, 금융 위기의 본질과 중앙은행의 역할에 초점을 맞춥니다. 우리는 1907년 공황, 대공황, 2008년 금융 위기 중 자동차 대출 시장 사례를 통해 금융 위기가 발생하는 원인과 그 파급 효과를 심층적으로 탐구합니다. 특히 중앙은행이 '최종 대부자(Lender of Last Resort)'로서 어떻게 금융 시장을 안정화하고 위기를 완화하는지, 그리고 그 과정에서 발생하는 정책적 한계와 비용은 무엇인지 살펴봅니다. 이 학습 자료는 비전공자도 금융 위기의 복잡한 메커니즘과 정책 대응을 명확하게 이해할 수 있도록 구성되었습니다.

\newpage

\section*{용어 정리}
금융 위기 관련 핵심 용어들을 쉽게 이해할 수 있도록 정리했습니다.

\begin{adjustbox}{width=\textwidth,center}
\begin{tabular}{lllc}
\toprule
\textbf{용어} & \textbf{쉬운 설명} & \textbf{원어} & \textbf{비고} \\
\midrule
뱅크런 & 많은 예금자가 동시에 은행에서 돈을 인출하려는 현상 & Bank Run & 은행 파산의 주요 원인 \\
만기 불일치 & 은행이 단기로 돈을 빌려(예금) 장기로 빌려주는(대출) 구조 & Maturity Mismatch & 은행 사업 모델의 핵심이자 취약점 \\
유동성 부족 & 은행에 당장 현금이 부족하지만, 자산 가치는 충분한 상태 & Illiquid & 일시적인 현금 부족 \\
지급 불능 & 은행의 자산 가치가 부채보다 적어 파산 상태인 경우 & Insolvent & 자본 잠식 상태 \\
숏 스퀴즈 & 주식을 빌려 판(공매도) 투자자들이 주가 상승으로 손실을 줄이기 위해 & Short Squeeze & 주가 급등을 유발하는 현상 \\
& 다시 주식을 사들이면서 주가가 더 오르는 현상 & & \\
최종 대부자 & 금융 시스템 위기 시, 다른 곳에서 돈을 빌릴 수 없을 때 & Lender of Last Resort & 중앙은행의 핵심 기능 \\
& 마지막으로 돈을 빌려주는 기관 (주로 중앙은행) & & \\
금본위제 & 통화 가치를 금에 고정하여, 화폐를 금으로 교환할 수 있도록 하는 제도 & Gold Standard & 통화량 조절에 제약이 있음 \\
자산유동화기업어음 & 자동차 대출 같은 자산을 담보로 발행하는 단기 어음 & Asset-Backed Commercial Paper (ABCP) & 비은행 금융기관의 주요 자금 조달원 \\
캡티브 금융 회사 & 자동차 제조사 등이 소유한 금융 회사로, 자사 제품 구매자에게 & Captive Finance Company & 자동차 대출 시장의 주요 플레이어 \\
& 대출을 제공함 & & \\
할인 창구 & 중앙은행이 시중 은행에 단기 자금을 대출해주는 제도 & Discount Window & 전통적인 중앙은행의 유동성 공급 수단 \\
프라이머리 딜러 & 중앙은행의 공개 시장 운영 거래 상대방이 되는 주요 투자은행 & Primary Dealer & 금융 시장의 핵심 중개자 \\
재정 정책 & 정부가 세금, 지출 등을 통해 경제에 영향을 미치는 정책 & Fiscal Policy & 정부 주도의 경제 안정화 정책 \\
통화 정책 & 중앙은행이 금리, 통화량 등을 조절하여 경제에 영향을 미치는 정책 & Monetary Policy & 중앙은행 주도의 경제 안정화 정책 \\
\bottomrule
\end{tabular}
\end{adjustbox}

\newpage

\section{모듈 4: 금융 위기}

\subsection{Lesson 4-0.1: 개요}
안녕하세요, 여러분. 이번 모듈에서는 중앙은행이 '최종 대부자(Lender of Last Resort)'로서 수행하는 역할에 집중합니다. 이를 위해 우리는 여러 금융 위기 사례를 깊이 있게 연구하고, 위기가 발생한 근본적인 원인을 논의할 것입니다. 이어서 중앙은행이 이러한 위기에 어떻게 정책적으로 대응했는지 살펴봅니다. 중앙은행의 노력이 성공을 거두기도 하지만, 단독으로 모든 문제를 해결하는 데는 한계가 있다는 점을 알게 될 것입니다.

\begin{tcolorbox}[title=\textbf{모듈 4 학습 목표}]
\begin{itemize}
    \item 다양한 금융 위기 사례를 통해 위기 발생 원인을 이해합니다.
    \item 중앙은행의 '최종 대부자' 역할과 그 중요성을 파악합니다.
    \item 금융 위기에 대한 중앙은행의 정책 대응 방안을 학습합니다.
    \item 중앙은행 단독 정책의 한계와 재정 정책의 필요성을 인식합니다.
\end{itemize}
\end{tcolorbox}

우리가 처음 살펴볼 위기는 1907년 공황(Panic of 1907)입니다. 이 위기는 20세기 최초의 주요 금융 위기였으며, 연방준비제도(Federal Reserve system)가 설립되기 전에 발생한 마지막 미국 은행 공황이기도 합니다. 이 사례를 통해 최종 대부자로서의 중앙은행이 왜 필요한지 명확하게 이해할 수 있습니다.

다음으로 대공황(Great Depression)을 연구합니다. 대공황은 오늘날 존재하는 사회 안전망(social safety net)의 대부분이 당시에는 없었다는 점에서 또 다른 유용한 사례입니다. 이를 통해 정부 개입이 없을 경우 경제가 어떻게 될 수 있는지 배울 수 있습니다.

그다음에는 금융 위기에 대한 국제적인 관점을 취하고, 위기의 원인을 파악하는 데 도움이 되는 공통된 주제들을 식별합니다. 이 거시적인 관점은 단기 자금이나 취약한 자금 조달(phone funding)로 조달되고 장기 프로젝트나 지속 불가능한 정부 정책에 투자된 '신용 붐(credit booms)'이 금융 위기를 예고했음을 강조합니다.

2008년 금융 위기 당시 미국 자동차 대출 시장을 연구함으로써 단기 자금 조달에 의존하는 위험을 다시 한번 강조합니다. 이 시기에 비은행 대출 기관들은 단기 자금 시장에 크게 의존하여 자동차 대출을 제공했습니다. 단기 자금 시장이 고갈되자 자동차 판매가 급감했습니다.

이 모든 사례를 살펴본 후, 여러분은 신용 붐이 항상 금융 위기와 경기 침체로 이어지는지 궁금해할 수 있습니다. 짧은 답변은 '아니오'입니다. 다음 강의에서 논의하겠지만, 이는 신용 붐이 어떤 부문에서 발생하는지에 따라 달라집니다. 예를 들어, 제조업 부문의 신용 붐은 적어도 역사적으로는 종종 지속 가능했습니다.

금융 위기의 원인과 결과를 검토한 후, 우리는 정책 대응으로 넘어갑니다. 구체적으로, 중앙은행이 뱅크런과 신용 공급의 급격한 감소에 어떻게 대응할 수 있는지 살펴봅니다. 이를 위해 2008년 금융 위기와 COVID-19 팬데믹에 대한 연방준비제도의 대응을 연구하여 '최종 대부자' 개념을 설명합니다. 즉, 연방준비제도는 다른 누구도 대출하지 않을 때 대출합니다.

통화 정책만이 금융 위기에 대응하는 유일한 방법은 아닙니다. 정부 이전(government transfers)과 지출을 포함하는 재정 정책(Fiscal Policy)도 경제를 도울 수 있습니다. 따라서 마지막 강의에서는 금융 위기에 대한 재정적 대응을 논의할 것입니다.

\newpage

\section{Lesson 4-1: 금융 위기 이해하기}

\subsection{Lesson 4-1.1: 뱅크런과 1907년 공황}
뱅크런(Bank Run)과 1907년 공황(Panic of 1907)에 대한 수업에 오신 것을 환영합니다. 우리는 두 단계로 진행할 것입니다. 첫째, 뱅크런이 무엇이며 은행이 왜 뱅크런에 취약한지 이해할 것입니다. 둘째, 우리의 첫 번째 주요 사례 연구인 1907년 공황을 살펴볼 것입니다. 이것은 미국 역사상 가장 두드러진 금융 위기 중 하나였습니다. 대규모 은행 공황이었으며, 현대 연방준비제도 시스템의 형성을 이끌었습니다. 우리는 공황이 왜 발생했으며 어떻게 해결되었는지 이해할 것입니다. 이러한 역사적 사건을 이해하는 것은 즐겁고 매우 유익할 수 있습니다. 정확한 역사적 세부 사항에 대해서는 너무 걱정하지 마십시오. 많은 세부 사항을 언급하겠지만, 이는 주로 재미를 위한 것이며, 서술과 우리가 배울 수 있는 교훈에 집중하십시오.

\subsubsection*{은행은 왜 뱅크런에 취약한가?}
은행은 단기적으로 예금을 받아 돈을 빌리고, 장기적으로 대출을 해줍니다. 이것은 자산(대출)과 부채(예금) 사이에 '만기 불일치(Maturity Mismatch)'를 만듭니다. 이러한 만기 변환(maturity transformation)은 은행 사업 모델의 핵심 특징입니다.

\begin{tcolorbox}[title=\textbf{핵심 개념: 만기 불일치 (Maturity Mismatch)}]
\textbf{한 줄 요약:} 은행이 단기 자금(예금)을 조달하여 장기 자산(대출)에 투자하는 구조입니다.

\textbf{직관적 비유:} 친구에게 "내일 갚을게" 하고 돈을 빌려, 그 돈으로 1년 뒤에나 수익이 나는 사업에 투자하는 것과 같습니다. 친구가 갑자기 "지금 당장 돈 갚아!"라고 하면 곤란해지겠죠?

\textbf{기술적 설명:} 예금자들은 언제든지 자신의 자금을 인출할 수 있지만, 은행이 대출해준 돈은 장기간 묶여 있습니다. 이 구조는 은행이 정상적으로 운영될 때는 효율적이지만, 예금자들이 동시에 돈을 인출하려 할 때 은행이 현금 부족에 직면하게 만듭니다.
\end{tcolorbox}

그렇다면 뱅크런은 왜 발생할까요? 예금자들은 언제든지 자신의 자금을 인출할 수 있다는 점을 기억하십시오. 어떤 충격이나 나쁜 소식이 발생하면, 이는 예금자들에게 위험을 초래합니다. 그들은 안전하게 돈을 인출하기로 결정할 수 있습니다. 아마도 은행 앞에 사람들이 줄 서 있는 것을 보았기 때문일 수도 있습니다. 뱅크런에 동참하는 것은 보통 합리적인 행동입니다. 만약 예금자가 줄에 합류하면, 돈을 돌려받을 기회가 있습니다. 줄의 맨 앞에 있는 사람들은 보통 돈을 돌려받을 것입니다. 만약 그들이 뱅크런에 동참하지 않으면, 돈을 잃을 수도 있습니다. 줄의 맨 뒤에 있는 사람들은 아무것도 받지 못할 수도 있습니다.

\subsubsection*{뱅크런의 결과}
많은 예금자가 동시에 도착하면, 은행은 충분한 현금을 가지고 있지 않을 수 있습니다. 대부분의 예금은 장기 투자에 묶여 있습니다. 이것이 만기 변환이 작동하는 방식입니다. 이 경우 은행은 '유동성 부족(illiquid)' 상태가 됩니다. 현금이 바닥난 것입니다.

\begin{tcolorbox}[title=\textbf{핵심 개념: 유동성 부족(Illiquid) vs. 지급 불능(Insolvent)}]
\textbf{한 줄 요약:} 유동성 부족은 현금이 일시적으로 없는 상태이고, 지급 불능은 자산보다 빚이 더 많은 파산 상태입니다.

\textbf{직관적 비유:}
\begin{itemize}
    \item \textbf{유동성 부족:} 지갑에 현금은 없지만, 비싼 집이나 차 같은 자산은 충분히 가지고 있는 상태입니다. 당장 현금이 필요할 때 자산을 팔아야 하지만, 팔 시간이 없거나 제값을 받기 어렵습니다.
    \item \textbf{지급 불능:} 지갑에 현금도 없고, 가지고 있는 모든 자산을 팔아도 빚을 다 갚을 수 없는 상태입니다. 완전히 파산한 것입니다.
\end{itemize}

\textbf{기술적 설명:}
\begin{itemize}
    \item \textbf{유동성 부족 (Illiquid):} 은행은 현금이 부족하지만, 그 투자는 여전히 수익성이 있을 수 있습니다. 유동성 부족 은행은 여전히 사업을 영위할 수 있으며, 지급 능력이 있을 수 있습니다.
    \item \textbf{지급 불능 (Insolvent):} 은행은 자본(equity)이 바닥나 파산 상태입니다.
\end{itemize}
\end{tcolorbox}

가장 최악의 시나리오는 건전한 은행이 뱅크런 때문에 파산하는 경우입니다. 이는 은행이 예금 인출 요구를 충족하기 위해 자산을 대폭 할인하여 강제로 매각해야 할 때 발생할 수 있습니다. 우리의 사례 연구에서 이러한 역학 관계가 실제로 어떻게 작용하는지 보게 될 것입니다. 그리고 이러한 최악의 시나리오를 방지하기 위해 다양한 금융 규제와 정부 기관이 존재하는 이유를 알게 될 것입니다.

\subsubsection*{사례 연구: 1907년 공황}
이제 우리의 사례 연구인 1907년 공황을 살펴보겠습니다. 1907년 공황은 뱅크런에 대해 배우는 데 유용한 사건입니다. 당시 현대 금융 시스템의 두 가지 특징이 없었습니다. 첫째, 중앙은행이 없었습니다. 이는 은행들이 정부 지원을 받는 비상 대출에 접근할 수 없었다는 것을 의미합니다. 둘째, 예금 보험이 없었습니다. 예금 보험이 없었기 때문에 예금자들은 훨씬 더 공황 상태에 빠졌습니다. 은행이 파산하면 평생 모은 돈을 잃을 수도 있었기 때문입니다. 이 역사적 사건은 뱅크런을 방지하기 위해 고안된 현대 정부 개입이 없을 때 어떤 일이 일어날 수 있는지 이해하는 데 도움이 됩니다. 이 금융 위기가 어떻게 전개되었는지 살펴보겠습니다.

대부분의 위기와 마찬가지로, 단일 원인이 아니라 나쁜 소식이 쌓여 뱅크런으로 이어졌습니다. 이야기는 1906년 샌프란시스코 지진으로 시작됩니다. 이 지진은 GDP의 약 1.5\%에 해당하는 피해를 입혔습니다. 이러한 대규모 충격은 즉시 금융 시스템 전체에 파급되었습니다. 보험사들은 보험금 청구를 충족하기 위해 자산을 매각해야 했습니다. 예금자들은 은행에서 돈을 인출했고, 뉴욕과 런던의 주식 시장은 하락했습니다. 이러한 추세는 계속되었습니다. 1907년 3월 며칠 동안 주식은 10\% 이상 하락했습니다. 단기 자금 시장(money markets)은 스트레스를 받았고, 이자율은 며칠 동안 2.5\%에서 10\% 이상으로 치솟았습니다. 이 사건은 단 하루 만에 일어난 것이 아니라 며칠에 걸쳐 발생했기 때문에 "조용한 폭락(the silent crash)"이라고 불리기도 합니다.

두 번째 충격은 1907년 작황이 전년도만큼 풍성하지 않아 수출이 감소한 것이었습니다. 미국 재무부는 만기가 도래하는 채권을 재융자할 수 없게 되자 은행에서 3천만 달러를 인출했습니다. 정부와 예금자들이 은행에서 더 많은 돈을 인출하면서 통화 공급이 위축되고 신용이 부족해지며 이자율이 더욱 상승했습니다. 요컨대, 1907년 여름 동안 금융 상황은 상당히 악화되었습니다.

이러한 배경 속에서 한 가지 더 중요한 요인이 있었습니다. 단기 은행 대출로 자금을 조달하여 주식 시장에서 위험한 투기를 하던 투기꾼들이었습니다. 이 위험한 투기가 실패하자 은행 대출이 부도 처리되었고, 이것이 은행 공황의 시작이었습니다. 이것이 짧은 버전의 이야기입니다. 여기 여러분이 흥미롭게 여길 만한 더 긴 버전이 있습니다.

하인츠 형제(Heinze process)는 매우 부유한 구리 광산 소유주였습니다. 1907년, 그들은 뉴욕시로 이주하여 찰스 W. 모스(Charles W. Morse)와 합류하여 구리 주식 시장에서 투기했습니다. 그들은 유나이티드 코퍼 컴퍼니(United Copper Company)의 많은 주식을 매입했습니다. 그들은 이 주식들을 담보로 삼아 20개에 달하는 은행들과 거래했습니다. 하인츠 형제와 그들의 동료들은 단기 은행 대출로 자금을 조달하여 주식 시장에서 위험한 투기를 하던 투기꾼들이었습니다.

1907년 10월 9일, 오토 하인츠(Otto Heinze)는 유나이티드 코퍼 주식의 거래 행태를 분석하기로 결정했습니다. 그는 이 주식 약 45만 주가 거래되었음을 발견했습니다. 그러나 회사에 따르면 발행 주식은 42만 5천 주에 불과했습니다. 이는 일부 브로커들이 주식을 빌려 공매도(shorted the stock)했다는 것을 의미했습니다. 즉, 이 다른 투자자들은 주식을 빌려 나중에 더 싼 가격에 다시 사들여 이익을 얻으려는 희망으로 주식을 팔았습니다. 하인츠 형제는 브로커들이 자신들이 소유하고 담보로 제공했던 주식을 빌려주었다고 믿었습니다. 그들은 주식을 회수함으로써 공매도자들이 자신들이 소유한 주식을 매우 높은 가격에 강제로 사들이게 되는 상황을 만들 수 있다고 판단했습니다. 이것을 '숏 스퀴즈(short squeeze)'라고 합니다.

\begin{tcolorbox}[title=\textbf{핵심 개념: 숏 스퀴즈 (Short Squeeze)}]
\textbf{한 줄 요약:} 공매도한 주식의 가격이 예상과 달리 급등하여, 공매도 투자자들이 손실을 줄이기 위해 주식을 다시 사들이면서 주가가 더욱 폭등하는 현상입니다.

\textbf{직관적 비유:} 친구에게 "나중에 싸게 사서 돌려줄게" 하고 연필을 빌려 팔았는데, 갑자기 연필 가격이 폭등한 상황입니다. 친구가 "지금 당장 연필 돌려줘!"라고 하면, 비싸게라도 연필을 사서 돌려줘야 하므로 손해가 커집니다. 이때 많은 공매도자들이 동시에 연필을 사들이면 가격은 더 오르게 됩니다.

\textbf{기술적 설명:} 하인츠 형제는 자신들이 담보로 제공한 주식을 회수하여 시장에 유통되는 주식 수를 줄임으로써, 공매도자들이 주식을 되사기 어렵게 만들고 가격을 올리려 했습니다.
\end{tcolorbox}

숏 스퀴즈를 유도하기 위해 하인츠 형제는 주식을 담보로 한 대출금을 갚기 위해 150만 달러를 빌려야 했습니다. 유나이티드 코퍼 주식이 45.50달러에서 37.75달러로 다시 하락하자, 하인츠 형제는 대출금 상환을 요청하며 주식을 회수하기 시작했습니다. 몇 시간 만에 주가는 23달러 상승했지만, 이후 약 52달러로 하락했습니다. 그러나 하인츠 형제는 오판했습니다. 주식을 담보로 대출을 연장했던 20개 은행 모두가 모든 주식을 인도했습니다. 결과적으로 하인츠 형제는 필요한 모든 대출금을 상환할 수 없었고, 주식 인수를 거부했습니다. 대출 기관들은 즉시 거부된 담보를 매각했고 주가는 폭락했습니다. 3일 후, 유나이티드 코퍼 주식은 가치의 50\%를 잃었고 결국 주당 10달러에 거래되었습니다. 숏 스퀴즈 계획은 처참하게 실패했습니다. 하인츠 형제는 파산했습니다.

10월 16일, 그들의 증권 중개 회사가 문을 닫았습니다. 하인츠 형제 중 한 명이 소유한 작은 은행도 마찬가지였습니다. 그러나 그 작은 은행은 뉴욕에 기반을 둔 머캔타일 은행(Mercantile Bank)을 결제 서비스에 이용했으며, 약 100만 달러의 예금을 가지고 있었습니다. 머캔타일 은행 또한 가치가 폭락한 유나이티드 코퍼 주식에 직접적으로 노출되어 있었습니다. 이러한 사건들은 너무나 충격적이어서 사람들은 하인츠 형제와 관련된 모든 은행에서 돈을 인출하기 시작했습니다. 공황이 시작된 것입니다.

모든 채권의 질서 있는 청산을 조정하기 위해 뉴욕 청산소(New York Clearing House), 즉 뉴욕 은행들의 컨소시엄이 10월 20일에 개입했습니다. 이것은 은행 공황을 진정시키기 위해 고안된 민간 부문의 개입이었습니다. 그러나 청산소 회의 중에 뉴욕에서 세 번째로 큰 은행인 니커보커 트러스트(Knickerbocker Trust)가 하인츠 형제와의 연관성 때문에 금융 부정에 연루되었다는 사실이 밝혀졌습니다. 이는 니커보커 트러스트가 심각한 문제에 처해 있음을 시사했습니다. 니커보커 트러스트가 재정적으로 건전했는지는 확실하지 않지만, 은행은 신뢰에 의존합니다. 예금자들이 돈을 돌려받을 수 있을지 의문을 품는다면, 그들은 단순히 돈을 인출할 것입니다.

그래서 1907년 10월 22일 아침, 니커보커 트러스트 사무실 문이 열리기 한 시간 전부터 줄이 형성되기 시작했습니다. 줄은 분 단위로 늘어나는 것 같았고, 모든 예금자들이 돈을 돌려달라고 요구했습니다. 줄 형성을 돕기 위해 경찰이 불려왔습니다. 정오까지 니커보커는 약 6,900만 달러의 자산과 6,400만 달러의 부채를 가지고 있으며 모든 사람에게 돈을 지급할 수 있다고 발표했습니다. 아무도 줄에서 이탈하지 않았습니다. 니커보커는 그날 800만 달러의 예금을 잃었습니다. 이것은 전형적인 뱅크런의 예입니다. 니커보커는 재정적으로 건전했지만, 걱정하는 예금자들이 어쨌든 돈을 돌려달라고 요구했습니다.

다음 날, 니커보커는 영업을 중단했습니다. 그러나 뉴욕의 다른 많은 건전한 은행들도 공황에 휩쓸렸습니다. 1907년 10월 24일까지 여러 은행이 파산했습니다. 트러스트 컴퍼니 오브 아메리카(Trust Company of America)에 대한 뱅크런이 시작되자, 뉴욕의 가장 유명한 은행가인 J.P. 모건(J.P. Morgan)과 다른 사람들이 개입했습니다. 그들은 예금자들의 요구를 충족시키기 위해 현금 대출을 제공했습니다. 현금 부족은 뉴욕 증권 거래소를 붕괴 직전으로 몰고 갔습니다. 다시 한번 J.P. 모건이 개입하여 증권 거래소를 유지하는 데 필요한 자본을 조달했습니다.

더 많은 은행들이 붕괴 직전에 처하자, J.P. 모건은 11월 초 뉴욕에 있는 자신의 집으로 약 120명의 뉴욕 은행장들을 소환했습니다. 그는 그들에게 위기 해결을 도울 것을 요구했습니다. 해결책을 강요하기 위해 그는 은행가들을 자신의 서재에 가두었다는 유명한 일화가 있습니다. 결국 은행가들은 동의했고, J.P. 모건은 위기 해결을 위해 2,500만 달러를 추가로 기부했습니다. 이 좋은 소식은 은행 공황을 멈추었지만, 이미 많은 피해가 발생했습니다. 은행 파산을 막기에는 너무 늦었고, 심각한 경기 침체를 막기에도 너무 늦었습니다.

\subsubsection*{1907년 공황의 교훈}
1907년 공황은 미국 역사상 가장 심각한 금융 위기 중 하나로 간주됩니다. 이 사건에서 무엇을 배울 수 있을까요?

\begin{tcolorbox}[title=\textbf{1907년 공황의 주요 교훈}]
\begin{enumerate}
    \item \textbf{신뢰 부족과 현금 부족:} 예금자들은 저축을 잃을까 봐 걱정했기 때문에 공황 상태에 빠졌습니다. 신뢰와 현금의 부족은 예금자들의 공황을 초래할 수 있습니다. 이것이 예금 보험이 뱅크런을 방지하는 이유를 설명합니다. 보험에 가입된 예금자들은 잠재적인 손실로부터 보호받을 것입니다. 이것이 오늘날 예금자들이 줄을 서는 고전적인 뱅크런이 드문 이유일 것입니다.
    \item \textbf{현금 대출의 중요성:} 두 번째 교훈은 현금 대출이 진행 중인 뱅크런을 예금자들의 요구를 충족시킴으로써 막는 데 도움이 될 수 있다는 것입니다. 1907년 공황에서 J.P. 모건은 '최종 대부자(Lender of Last Resort)' 역할을 했습니다.
    \item \textbf{정부 지원의 한계:} 1907년에는 J.P. 모건이 현금이 바닥나거나 자신의 은행에 문제가 생기지 않았다는 점에서 운이 좋았습니다. 정부의 지원 없이는 민간 최종 대부자의 선택지는 제한적입니다. 이것이 무제한적인 자원을 가진 연방준비제도가 최종 대부자로서 더 효과적일 수 있는 이유를 설명합니다.
\end{enumerate}
\end{tcolorbox}

\newpage

\subsection{Lesson 4-1.2: 대공황}
대공황(The Great Depression)에 대한 수업에 오신 것을 환영합니다. 현대 시대에 대공황은 미국과 서구 세계 대부분에서 가장 심각한 금융 위기이자 경제 침체였습니다. 우리는 대공황 동안 금융 상황이 어떻게 전개되었는지 살펴볼 것입니다. 우리는 엄청난 수의 은행 파산을 보게 될 것입니다. 우리는 연방준비제도 시스템이 은행 공황을 막는 데 왜 비효율적이었는지 이해할 것입니다. 마지막으로, 경제를 재활성화하고 은행 부문에 대한 신뢰를 회복시킨 대규모 정부 대응을 살펴볼 것입니다. 여기에는 정부 지출 프로그램이 포함되었습니다. 또한 예금 보험과 같은 새로운 기관의 설립도 포함되었습니다. 이러한 역사적 사건을 이해하는 것은 즐겁고 매우 유익할 수 있습니다. 정확한 역사적 세부 사항에 대해서는 너무 걱정하지 마십시오. 많은 세부 사항을 언급하겠지만, 이는 주로 재미를 위한 것이며, 서술과 우리가 배울 수 있는 교훈에 집중하십시오. 대공황은 예금 보험이나 실업 수당과 같이 금융 위기의 영향을 완화할 대부분의 정책이 미국에 존재하지 않았다는 점에서 흥미롭습니다. 이는 정부가 금융 위기 동안 개입하지 않았을 때 경제에서 어떤 일이 일어날지 이해하는 데 도움이 됩니다.

\subsubsection*{대공황의 시작}
대공황의 시작을 간략하게 요약해 보겠습니다. 1929년 10월 24일, 이른바 '검은 목요일(Black Thursday)'에 개장과 동시에 시장은 11\% 하락했습니다. 이 폭락은 대규모 투자자들에 대한 사기 혐의로 인한 런던 증권 거래소의 하락에 뒤이어 발생했습니다. 1929년 10월 28일, 주가 하락으로 인해 점점 더 많은 투자자들이 신용으로 매수한 주식에 대한 마진콜(margin calls)에 직면했습니다. 결과적으로 더 많은 투자자들이 주식을 팔아 손실을 줄였습니다. 다우 지수는 그날 거의 13\% 하락했습니다. 다음 날, 투자자들은 기록적인 수의 주식을 거래했습니다. 하루 만에 1,600만 주가 거래되었습니다. 다음 날, 주식 시장은 유동성 부족의 징후를 보였습니다. 일부 주식은 어떤 가격에도 팔리지 않았습니다. 다우 지수는 또 다른 12\% 하락했습니다.

\begin{tcolorbox}[title=\textbf{주식 시장 붕괴 (1929년 9월 ~ 1932년 7월)}]
\textbf{핵심 내용:}
\begin{itemize}
    \item 1929년 10월의 사건들은 주가 장기 하락의 서막에 불과했습니다.
    \item 1929년 9월부터 1932년 7월까지 주가는 80\% 이상 폭락했습니다.
\end{itemize}
\textbf{시사점:} 주식 시장의 붕괴는 단순한 일시적 현상이 아니라, 장기적인 경제 침체의 시작을 알리는 신호였습니다.
\end{tcolorbox}

위기가 시작될 당시, 8,000개 이상의 상업 은행이 연방준비제도 시스템에 속해 있었지만, 거의 16,000개는 그렇지 않았습니다. 당시 은행들은 종종 주식을 직접 소유하고 있었고 현금 준비금이 제한적이어서, 연방준비제도 시스템에 접근하여 현금을 확보할 수 없다는 것은 문제가 되었습니다. 연방준비제도에 접근할 수 없는 은행들은 종종 다른 대형 은행에 현금을 예치했는데, 이는 소규모 은행이 대형 은행에서 예금을 인출하거나 대형 은행 자체가 어려움에 처할 경우 전염(contagion) 가능성을 만들었습니다. 이러한 환경에서 예금자들은 당연히 은행의 건전성에 대해 의문을 품었습니다. 대공황의 어쩌면 놀라운 특징은 하나의 전국적인 은행 공황이 아니라 수많은 지역 은행 뱅크런이 있었다는 것입니다. 총 1929년에서 1933년 사이에 약 9,000개의 은행이 문을 닫았습니다.

\begin{tcolorbox}[title=\textbf{은행 파산의 파급 효과}]
\textbf{핵심 내용:}
\begin{itemize}
    \item 대규모 은행 폐쇄는 예금을 급격히 감소시켰습니다.
    \item 은행들은 현금 준비금을 확보하기 위해 대출을 줄였습니다.
    \item 결과적으로 통화 공급이 위축되었습니다.
\end{itemize}
\textbf{시사점:} 은행 파산은 단순한 개별 은행의 문제가 아니라, 전체 금융 시스템의 신용 경색과 통화 공급 감소로 이어져 경제 전반에 악영향을 미쳤습니다.
\end{tcolorbox}

여러분은 연방준비제도가 왜 통화 공급을 늘리지 않았는지 궁금할 수 있습니다. 당시 미국은 다른 대부분의 국가와 마찬가지로 금본위제(gold standard)를 채택하고 있었습니다. 이는 화폐가 고정된 환율로 금에 연동되어 있었고, 이는 중앙은행의 선택지를 심각하게 제한했습니다.

\subsubsection*{디플레이션의 심화}
낮은 통화 공급은 물가에 하방 압력을 가했습니다. 이 차트에서 볼 수 있듯이, 대공황 동안 소비자 물가는 약 25\% 디플레이션되었습니다. 식품의 경우 하락폭이 훨씬 더 컸습니다. 이러한 디플레이션은 농부들과 기업들이 실질적으로, 즉 생산된 상품의 측면에서 더 높은 부채에 직면했음을 의미했습니다. 결과적으로 많은 농부들과 기업들이 파산했습니다. 이러한 파산은 은행 시스템에 추가적인 스트레스를 유발했습니다.

\begin{tcolorbox}[title=\textbf{핵심 개념: 디플레이션 (Deflation)}]
\textbf{한 줄 요약:} 물가가 지속적으로 하락하는 현상입니다.

\textbf{직관적 비유:} 오늘 100원짜리 빵이 내일 90원이 되고 모레 80원이 된다고 상상해 보세요. 사람들은 "내일 사면 더 싸지겠네?" 하고 소비를 미루게 됩니다. 기업은 물건이 안 팔리니 생산을 줄이고, 직원도 해고하게 됩니다.

\textbf{기술적 설명:} 물가 하락, 특히 미래의 물가 하락 예상은 투자를 할 유인을 감소시킵니다. 만약 현금을 보유하면 미래에 더 많은 상품을 살 수 있다면, 왜 미래에 낮은 가격만 청구할 수 있는 위험한 사업에 현금을 투자하겠습니까? 이는 경제 활동을 더욱 위축시킵니다.
\end{tcolorbox}

디플레이션이 시작되고 신용 가용성이 줄어들면서 기업 투자는 급락했습니다. 투자 감소는 대규모 파산과 함께 산업 생산의 장기적인 감소로 이어졌습니다. 이 차트에서 볼 수 있듯이, 대공황 동안 산업 생산량은 50\% 감소했습니다. 이는 전례 없는 경제 활동의 감소였습니다.

\begin{tcolorbox}[title=\textbf{경제 활동 위축 (산업 생산 및 실업률)}]
\textbf{핵심 내용:}
\begin{itemize}
    \item 산업 생산량은 대공황 동안 50\% 급감했습니다.
    \item 실업률은 대공황 초기의 거의 0\%에서 끝 무렵에는 25\% 이상으로 증가했습니다.
\end{itemize}
\textbf{시사점:} 경제가 위축되면서 일자리가 급격히 사라졌고, 이는 빈곤과 사회 불안으로 이어졌습니다. 실업 보험이 없었기 때문에 실업자들은 더욱 큰 고통을 겪었습니다.
\end{tcolorbox}

경제 위축과 함께 실업률이 빠르게 상승했습니다. 이 차트에서 볼 수 있듯이, 실업률은 대공황 초기의 거의 0\%에서 대공황 말기에는 25\% 이상으로 증가했습니다. 실업 보험이 존재하지 않았기 때문에 빈곤과 사회 불안이 뒤따랐습니다.

\subsubsection*{사회 불안과 정부 대응}
사회 불안의 가장 두드러진 예 중 하나는 1차 세계대전 참전 용사들의 행진이었습니다. 1932년, 1차 세계대전 참전 용사들은 '보너스 원정군(Bonus Expeditionary Forces)'이라는 단체를 조직하여 워싱턴 D.C.로 행진했습니다. 약 2만 명의 참전 용사들이 수도에 모여 워싱턴 D.C.에 야영하며 약속되었지만 의회에서 지연되고 있던 보너스 지급을 요구했습니다. 시위는 결국 강제로 해산되었지만, 굶주린 참전 용사들을 해산시키는 탱크의 사진은 후버 대통령의 재선 도전을 약화시켰습니다. 11월, 프랭클린 델라노 루스벨트(Franklin Delano Roosevelt)는 경제 포퓰리즘을 내세워 선거에서 승리했고, 참전 용사 보너스는 결국 1936년에 지급되었습니다.

루스벨트가 취임한 후 첫 번째 조치 중 하나는 1933년 3월 6일에 '뱅크 홀리데이(bank holiday)', 더 정확히는 '뱅크 홀리데이 주간'을 선언한 것이었습니다. 일주일 내내 아무도 은행이나 은행 서비스에 접근할 수 없었습니다. 이는 돈을 인출하거나 이체할 수 없다는 것을 의미했으며, 그해 초 다시 시작되었던 뱅크런을 중단시켰습니다. 뱅크 홀리데이의 목적은 은행 시스템에 대한 신뢰를 회복하고 미래의 뱅크런을 방지할 법안을 제정할 시간을 버는 것이었습니다.

1933년 긴급 은행법(Emergency Banking Act of 1933)은 이 목표를 달성했습니다. 요컨대, 이 법은 건전한 은행만이 재개하고 부실한 은행은 폐쇄되도록 보장했습니다. 또한 이 법은 연방준비제도에 더 많은 권한을 부여하고 은행에 대출하기 쉽게 만들었습니다. 연방준비제도는 사실상 재개된 은행의 예금 보증인이 되었습니다.

\begin{tcolorbox}[title=\textbf{긴급 은행법의 즉각적인 효과}]
\textbf{핵심 내용:}
\begin{itemize}
    \item 법안 통과 후 다우존스 산업 주가 지수는 하루 만에 15\% 이상 상승했습니다.
    \item 뱅크런이 중단되었습니다.
\end{itemize}
\textbf{시사점:} 정부의 신속하고 단호한 조치가 시장의 신뢰를 회복하고 금융 공황을 진정시키는 데 결정적인 역할을 했음을 보여줍니다.
\end{tcolorbox}

1933년 은행법(Banking Act of 1933), 일명 글래스-스티걸 법(Glass-Steagall Act)은 1933년 6월에 예금자 보호를 강화했습니다. 상업 은행(예금 수취)과 투자 은행을 분리함으로써, 예금자들은 더 이상 주식이나 일반적으로 증권 시장에서의 은행의 위험 감수에 노출되지 않게 되었습니다. 이 법은 또한 연방예금보험공사(Federal Deposit Insurance Corporation, FDIC)를 설립하여 예금에 대한 보험을 제공함으로써 추가적인 뱅크런을 방지하는 데 도움이 되었습니다. 은행법은 또한 은행 감독을 강화하고 연방준비제도 시스템의 일부 특징을 변경했습니다. 은행 공황은 끝났지만, 실물 경제는 여전히 심각한 상황에 처해 있었습니다.

경제 위기에서 정부의 역할을 근본적으로 재편한 '뉴딜(New Deal)' 정책의 일환으로, 루스벨트의 첫 번째 목표는 실업자들의 고통을 완화하는 것이었습니다. 정부는 공공사업진흥국(Works Progress Administration, WPA)과 민간산림보존단(Civilian Conservation Corps, CCC)을 통해 일자리를 창출했습니다. 공공사업진흥국은 주로 건설 분야에서 약 850만 명에게 일자리를 제공했습니다. 65만 마일 이상의 도로, 12만 5천 개의 공공 건물, 7만 5천 개의 다리, 8천 개의 공원이 건설되었습니다. 민간산림보존단은 나무 심기, 홍수 방지벽 건설, 산불 진압, 산림 도로 유지 보수 등의 작업을 제공했습니다. 이러한 프로젝트들은 단기간에 시작될 수 있었고 즉각적인 효과를 가져왔다는 점이 중요합니다.

\begin{tcolorbox}[title=\textbf{뉴딜 정책의 고용 효과}]
\textbf{핵심 내용:}
\begin{itemize}
    \item 뉴딜 정책 시행 후 1년 이내에 실업률이 9\% 포인트 감소했습니다.
\end{itemize}
\textbf{시사점:} 정부의 대규모 공공 사업과 일자리 창출 프로그램이 실업 문제 해결에 직접적이고 효과적인 영향을 미쳤음을 보여줍니다.
\end{itemize}
\end{tcolorbox}

사람들의 경제적 어려움은 또한 사회 안전망의 필요성을 보여주었습니다. 사회보장국(Social Security Administration)을 통한 연금 시스템 구축, 실업 수당, 장애 수당은 '뉴딜'의 일부였습니다. '뉴딜'은 또한 전국노동관계위원회(National Labor Relations Board)를 설립하여 노동 관계를 재편했습니다. 경제는 다시 정상 궤도에 오르는 것처럼 보였습니다. 그러나 '뉴딜'은 자금 조달이 필요했고, 세금이 인상되었습니다. 더욱이 경제가 개선되면서 재정 부양책이 철회되었습니다. 동시에 연방준비제도는 은행의 지급준비율(reserve requirements)을 인상하여 신용 가용성을 감소시켰습니다.

\begin{tcolorbox}[title=\textbf{1937년 경기 침체: 정책 오류의 경고}]
\textbf{핵심 내용:}
\begin{itemize}
    \item 1937년, 긴축 통화 정책(연준의 지급준비율 인상)과 긴축 재정 정책(재정 부양책 철회)이 동시에 시행되었습니다.
    \item 결과적으로 산업 생산이 다시 붕괴되었습니다.
\item 경기 침체는 연방준비제도가 지급준비율을 인하하면서 끝났습니다.
\end{itemize}
\textbf{시사점:} 1937년 사건은 경제 지원을 너무 일찍 철회하는 것에 대한 경고를 제공합니다. 위기에서 벗어나기 위한 정책적 노력이 충분히 지속되지 않으면, 경제는 다시 침체에 빠질 수 있습니다.
\end{tcolorbox}

이 차트에서 볼 수 있듯이, 1937년 경기 침체 이후 실업률은 빠르게 감소했습니다. 그러나 또 다른 부양책이 경제를 빠르게 성장시키는 데 도움이 되었다는 점에 유의하는 것이 중요합니다. 유럽의 국방력 증강과 그에 따른 정부 지출 증가였습니다. 1940년 중반부터 실업률 감소가 가속화되는 것을 볼 수 있습니다.

\subsubsection*{대공황에서 얻은 교훈}
이 강의에서 무엇을 배웠습니까?

\begin{tcolorbox}[title=\textbf{대공황의 주요 교훈}]
\begin{enumerate}
    \item \textbf{금융 위기의 파괴적 영향:} 대공황을 금융 위기의 결과에 대한 사례 연구로 사용하면서, 우리는 그 영향이 특히 사회 안전망이나 정부 개입이 없을 때 얼마나 파괴적일 수 있는지 보았습니다.
    \item \textbf{금융 시스템의 대규모 변화 필요성:} 금융 시스템에 대한 신뢰를 회복하려면 대규모 변화가 필요합니다.
    \item \textbf{정부 지원의 조기 철회 위험:} 1937년의 경기 침체는 경제에 대한 정부 지원을 너무 적게 제공하거나 너무 일찍 철회하는 것이 또 다른 경제 침체로 이어질 수 있음을 보여줍니다.
\end{enumerate}
\end{tcolorbox}

\newpage

\subsection{Lesson 4-1.3: 금융 위기의 공통 특징}
금융 위기의 공통 특징에 대한 수업에 오신 것을 환영합니다. 우리는 지난 150년 동안의 주요 글로벌 금융 위기에 대한 광범위한 개요를 제공할 것입니다. 우리의 목표는 위기가 어떻게 전개되었는지, 왜 위기가 전 세계적이었는지, 그리고 이러한 위기들 사이에 어떤 유사점이 있는지 이해하는 것입니다. 우리는 금융 위기의 이러한 공통된 주제들을 밝혀낼 것입니다. 과도한 부채가 항상 역할을 한다는 것을 알게 될 것입니다. 은행 예금과 같은 단기 부채에 대한 런이 항상 있습니다. 지속 불가능한 장기 투자나 정부 정책도 종종 역할을 합니다. 신용의 급격한 확장이 항상 금융 위기로 끝나는가요? 강의가 끝날 무렵에 이 질문에도 답할 것입니다. 시작해 봅시다.

\subsubsection*{주요 글로벌 금융 위기 사례}
지난 150년 동안 가장 심각했던 5가지 글로벌 금융 위기는 1873년 공황(Panic of 1873), 1890년 베어링스 위기(Barings Crisis of 1890), 1907년 공황(Panic of 1907), 1930년대 대공황(Great Depression of the 1930s), 그리고 2008년 금융 위기(2008 Financial Crisis)였습니다. 각각을 차례로 살펴보겠습니다. 이것은 간략한 설명이며, 우리의 목표는 계속해서 나타나는 특징들을 찾는 것임을 기억하십시오. 모든 작은 세부 사항보다는 큰 그림에 집중하십시오.

\begin{tcolorbox}[title=\textbf{금융 위기의 공통 특징을 찾는 여정}]
우리는 다음의 주요 글로벌 금융 위기들을 통해 공통된 패턴을 찾아낼 것입니다.
\begin{itemize}
    \item \textbf{1873년 공황:} 철도 투자 버블과 국제적 확산
    \item \textbf{1890년 베어링스 위기:} 아르헨티나 정부 부채와 국제적 자금 조달 문제
    \item \textbf{1907년 공황:} 구리 투기와 은행 시스템의 취약성
    \item \textbf{1930년대 대공황:} 금본위제와 고정 환율을 통한 전 세계적 전파
    \item \textbf{2008년 금융 위기:} 주택 시장 붕괴와 글로벌 금융 시스템의 상호 연결성
\end{itemize}
\end{tcolorbox}

\paragraph{1873년 공황}
1873년 공황은 최초의 글로벌 은행 위기였습니다. 철도는 1800년대 중반의 뜨거운 투자 대상이었습니다. 1869년 미국 동부와 서부 해안을 연결하는 대륙횡단철도(Transcontinental railroad) 완성과 같은 이정표는 철도 투자에 대한 관심을 부추겼습니다. 이는 미국뿐만 아니라 유럽에서도 마찬가지였습니다. 예를 들어, 독일 전역에서 부채로 자금을 조달한 철도 회사들이 생겨났습니다. 철도 주가는 전 세계적으로 급등했습니다. 많은 은행들이 철도 채권을 투자로 보유했습니다. 한 국가의 은행 시스템의 주요 취약점은 철도 채권 투기였습니다. 실제로 철도는 많은 선행 투자가 필요했고 초기 몇 년 동안은 수익이 거의 또는 전혀 발생하지 않아 투자자들에게 위험을 안겨주었습니다.

철도 투기 거품 붕괴로 이어진 여러 요인들이 있었습니다. 예를 들어, 위기 직전 보스턴과 시카고의 화재는 은행 준비금을 감소시켰습니다. 독일과 미국은 은본위제(silver currencies)를 포기했는데, 이는 광부들에게 부담을 주고 철도 네트워크의 필요성을 감소시켰습니다. 그리고 프랑스-독일 전쟁으로 인한 혼란이 여전히 존재했습니다. 위기의 첫 징후는 5월 9일 비엔나에서 주식 시장이 폭락하고 일련의 파산이 뒤따르면서 나타났습니다. 독일에서는 대규모 철도 회사가 붕괴되어 공황을 가중시켰습니다. 미국에서는 주요 은행인 J. 코흐 앤 컴퍼니(J Koch and Company)가 철도 채권을 매각할 수 없게 되어 9월 중순에 붕괴되었고, 이는 미국 전역에서 뱅크런과 파산을 촉발했습니다.

미국 은행 위기의 저명한 희생자 중 하나는 프리드먼 저축 은행(Friedman Savings Bank)이었습니다. 이 은행은 당시 자유를 얻은 전 노예들을 위한 최초이자 유일한 은행이었습니다. 미국과 유럽은 실업률이 빠르게 상승하는 깊은 경기 침체에 빠졌습니다. 이 위기는 오스만 제국과 라틴 아메리카에도 영향을 미쳤습니다. 종합적으로 볼 때, 이 위기는 투기적 투자, 즉 철도 채권에서 시작되어 대규모 손실과 은행 시스템에 대한 신뢰 부족으로 이어졌습니다.

\paragraph{1890년 베어링스 위기}
정부 부채라는 다른 유형의 투자가 이른바 1890년 베어링스 위기를 촉발했습니다. 1890년까지 아르헨티나는 10년간 대규모 자본 유입을 경험했습니다. 자본은 주로 런던에서 베어링스 은행(Barings bank)을 포함한 유럽 은행들에 의해 조달되었습니다. 이 돈은 철도 및 운송 네트워크 건설, 농지 개선을 포함한 장기 인프라 투자 프로젝트에 자금을 조달하는 데 사용되었습니다. 그러나 1880년대 후반까지 해외 차입금의 40\%가 부채 상환에, 수입의 60\%가 소비재에 사용되었습니다. 아르헨티나 정부 또한 상당한 재정 적자를 겪고 있었습니다. 철도 순이익은 1888년에 정점을 찍었고, 결국 국제 투자자들은 아르헨티나 정부의 부채 부담에 대해 경계심을 갖게 되었습니다. 결국 베어링스는 아르헨티나의 채권 발행을 위한 투자자를 찾을 수 없었습니다. 아르헨티나는 뱅크런을 경험했고, 새로운 자금을 조달할 수 없게 되자 정부 채무 불이행(default)을 선언했습니다. 이 채무 불이행은 라틴 아메리카에 심각한 경기 침체를 초래했을 뿐만 아니라 유럽과 미국의 주식 시장에도 급격한 하락을 가져왔습니다. 1873년 위기와 마찬가지로, 투기적인 장기 투자에 대한 신용 붐은 결국 나쁜 결과를 초래했습니다. 그러나 국제적 연계를 주목하십시오. 투기 활동은 국제 투자자들에 의해 부추겨졌습니다.

\paragraph{1907년 공황}
다음 위기인 1907년 공황은 구리 투기에서 시작되었습니다. 1907년 공황은 별도의 강의에서 논의했으므로 간략하게 설명하겠습니다. 샌프란시스코 지진은 보험 회사, 철도, 은행에 큰 손실을 입혔습니다. 이러한 배경 속에서 구리 주식 투기가 잘못되어 은행 시스템에 더 많은 손실을 초래했습니다. 이는 은행 공황, 은행 폐쇄, 그리고 미국에서 유럽으로 확산된 경기 침체로 이어졌습니다. 미국에서는 실업률이 16\% 이상으로 상승했습니다. 글로벌 시장도 큰 타격을 입었습니다. 예를 들어, 영국의 실업률은 1906년과 1908년 사이에 3.6\%에서 7.8\%로 두 배가 되었습니다. 독일, 일본, 이탈리아도 뱅크런을 경험했습니다. 1907년 공황은 금융 시장 투기가 잘못된 예입니다. 그러나 그 영향은 다른 위기와 유사합니다. 은행 손실은 뱅크런을 유발했습니다. 그리고 은행 시스템의 혼란은 신용과 실물 경제를 위축시켰습니다. 여기에도 국제적 연계가 있었습니다. 여기에는 샌프란시스코의 재해 위험을 인수하는 해외 보험 회사와 외국인 투자자들의 위험한 주식 시장 투자가 포함되었습니다. 이는 위기가 미국에서 전 세계로 확산되도록 했습니다.

\paragraph{대공황}
대공황은 국제적 연계의 중요성을 가장 명확하게 보여주는 예입니다. 대공황은 별도로 논의했으므로 국제적 요소에만 집중하겠습니다. 1929년 미국에서 주식 시장 붕괴가 대공황을 촉발했을 때, 대부분의 국가들은 금본위제(gold standard)를 채택하고 있었습니다. 금본위제는 모든 지폐가 금으로 전환될 수 있었고, 따라서 금 준비금이 통화 공급을 결정한다는 것을 의미했습니다. 이는 또한 각 통화가 금에 고정되어 있었기 때문에 대부분의 국가가 고정 환율 시스템으로 운영되었다는 것을 의미했습니다.

당시 미국은 글로벌 시장에서 금의 순 대출국이었습니다. 그러나 은행 위기가 전개되자 미국은 금을 대출하기보다는 비축했습니다. 이는 다른 국가들이 각자의 통화 공급 감소를 피하기 위해 금 준비금을 유지하기 위해 고군분투하게 만들었습니다. 낮은 통화 공급은 은행에 부담을 주고 전 세계적으로 디플레이션을 초래했습니다. 예를 들어, 독일 라이히스방크(Reichsbank)는 단기 채권자들이 독일에 대한 대출을 중단했을 때 사실상 금이 바닥났습니다. 결국, 국가들이 자체 통화 공급을 통제하고 뱅크런과 디플레이션을 피할 수 있도록 금본위제는 폐지되었습니다. 요컨대, 대공황은 고정 환율 때문에 글로벌 금융 시스템을 통해 전파되었습니다. 따라서 이전 금융 위기와 달리, 여기서는 나쁜 장기 투자에 대한 공통된 노출이 문제가 아니었습니다. 이번에는 전 세계의 모든 통화와 모든 통화 공급이 연결되어 있었습니다. 변동 환율(floating exchange rates)에서는 국제적 충격에 직면했을 때 통화 가치가 조정됩니다. 고정 환율에서는 통화 공급이 조정되어야 하며, 이는 큰 부정적인 결과를 초래합니다.

\paragraph{2008년 금융 위기}
이제 2008년 금융 위기로 넘어갑니다. 이 사건이 과거 위기들과 얼마나 유사했는지 이제 알 수 있을 것입니다. 위기는 미국 주택 시장에서 시작되었습니다. 스페인과 아일랜드와 같은 다른 많은 국가들도 주택 붐과 붕괴를 경험했습니다. 주택 붐에 자금을 조달하는 많은 장기 모기지 투자는 단기 신용으로 조달되었습니다. 이 단기 신용이 고갈되자 모기지 시장과 주택 가격이 붕괴되었습니다. 이전 금융 위기와 마찬가지로, 국제 투자자들은 이러한 투자에 노출되어 있었습니다. 이러한 노출의 정도는 2007년 6월 중견 독일 은행인 ECA Bay Deutsche Industriebank가 미국 서브프라임 모기지 손실로 파산했을 때 명확해졌습니다. 또한 국내 플레이어들, 가장 악명 높은 투자 은행인 리먼 브라더스(Lehman Brothers)가 모기지로 담보된 단기 자금 조달에 의존하고 있다는 것이 분명해졌습니다.

리먼 브라더스는 2008년 가을에 붕괴되었습니다. 가장 큰 보험 회사인 아메리칸 인터내셔널 그룹(American International Group)도 마찬가지였습니다. 이는 전 세계적으로 주식 시장 폭락과 자금 시장 공황을 촉발했습니다.

2008년 위기는 자본이 국경을 자유롭게 넘나드는 글로벌 통합 금융 시스템 때문에 전 세계적으로 전개되었습니다. 글로벌 통합 금융 시스템은 충격이 작을 때 쉽게 흡수될 수 있도록 돕습니다. 그러나 가장 큰 경제에 대한 충격은 모든 사람에게 영향을 미칩니다.

\subsubsection*{신용 붐은 항상 금융 위기로 이어지는가?}
마무리하기 전에 마지막 질문을 살펴보겠습니다. 신용 붐이 항상 금융 위기로 이어지는가요? 짧은 답변은 '아니오'입니다. 최근 연구에 따르면 많은 신용 붐이 금융 위기로 이어지지 않고, 대신 지속적인 경제 성장을 가져올 수 있습니다. 이 차트는 1910년부터 2014년까지 189개국 표본에 대한 신용 붐 정점 전후의 생산량 역학을 보여줍니다. 신용 붐의 정점은 이벤트 시간 0, 즉 수직선으로 표시됩니다. 두 가지 다른 유형의 신용 붐을 나타내는 두 개의 선이 있습니다. 대략적으로 말하면, 하나는 건설 및 부동산 부문이고 다른 하나는 제조업 부문입니다.

\begin{tcolorbox}[title=\textbf{신용 붐과 경제 성장: 부동산 vs. 제조업}]
\textbf{차트 분석 (텍스트 기반 설명):}
\begin{itemize}
    \item \textbf{신용 붐 기간:} 평균적으로 GDP는 붐 기간 동안 강하게 성장합니다.
    \item \textbf{부동산 신용 붐 이후:} 정점 이후에만 생산량이 감소하는 경향을 보입니다.
    \item \textbf{제조업 신용 붐 이후:} 생산량이 지속적으로 성장하거나 안정적인 경향을 보입니다.
\end{itemize}
\textbf{왜 이런 현상이 발생할까요?}
\begin{itemize}
    \item \textbf{부동산 신용 붐:} 일반적으로 주택 가격 상승, 신규 주택 건설, 고용 증가 등을 유발합니다. 그러나 결국 붐을 지속할 만큼 충분한 지역 구매자가 없게 됩니다. 이는 주택 가격에 하방 압력을 가하고, 건설 부문의 붕괴로 이어질 수 있습니다.
    \item \textbf{제조업 신용 붐:} 글로벌 수요나 신제품 수요에 의해 주도될 때 더 지속 가능할 수 있습니다. 예를 들어, 1960년대부터 1980년대까지 한국 제조업의 빠른 성장을 생각해 보십시오. 대부분의 제품은 수출되었습니다. 제조업 부문에 대한 신용은 수익성 있는 투자에 자금을 조달할 수 있습니다. 이러한 더 높은 이익은 신용 붐 동안 발생한 부채를 상환하는 데 도움이 됩니다.
\end{itemize}
\textbf{결론:} 모든 신용 붐이 금융 위기로 이어지는 것은 아니며, 어떤 부문에서 신용 붐이 발생하는지가 중요합니다.
\end{tcolorbox}

\subsubsection*{금융 위기의 공통 특징 요약}
이 세션에서 우리는 금융 위기의 공통 특징을 설명했습니다.

\begin{tcolorbox}[title=\textbf{금융 위기의 공통 특징}]
\begin{itemize}
    \item \textbf{과도한 부채:} 특히 단기 부채가 많습니다.
    \item \textbf{실패한 장기 투자:} 결국 나쁜 결과를 초래하는 장기 투자가 있습니다.
    \item \textbf{지속 불가능한 정부 정책:} 과도한 지출이나 고정 환율과 같은 지속 불가능한 정부 정책이 위기를 심화시킵니다.
    \item \textbf{국제적 전파:} 일부 금융 위기는 국제적 연계를 통해 전 세계적으로 확산됩니다.
\end{itemize}
\textbf{추가 교훈:} 모든 신용 붐이 은행 위기나 GDP 하락으로 이어지는 것은 아닙니다. 역사적으로 제조업 부문의 신용 붐은 더 지속 가능했습니다.
\end{itemize}
\end{tcolorbox}

\newpage

\subsection{Lesson 4-1.4: 2008년 금융 위기 중 자동차 대출}
2008년 금융 위기 중 자동차 대출에 대한 강의에 오신 것을 환영합니다. 이 사례 연구에서는 단기 자금 조달이 고갈되었을 때 비은행 대출 기관이 큰 역할을 하는 자동차 대출 시장에서 어떤 일이 일어났는지 살펴볼 것입니다. 은행뿐만 아니라 다른 금융 기관도 런에 취약하며, 신용 공급의 감소가 자동차 판매 감소로 이어진다는 것을 알게 될 것입니다.

\subsubsection*{캡티브 금융 회사와 자금 조달}
2008년 금융 위기 이전에 전체 자동차 대출의 약 절반은 '캡티브 금융 회사(captive finance companies)'에 의해 시작되었습니다. 캡티브 금융 회사는 자동차 제조업체가 소유한 대출 기관으로, 예를 들어 포드 모터 크레딧(Ford Motor Credit), GM 파이낸셜(GM Financial), 토요타 파이낸셜 서비스(Toyota Financial Services) 등이 있습니다.

\begin{tcolorbox}[title=\textbf{캡티브 금융 회사의 시장 점유율}]
\textbf{핵심 내용:}
\begin{itemize}
    \item 캡티브 금융 회사의 시장 점유율은 자동차 제조업체가 위치한 미시간주, 그리고 역사적으로 북부보다 은행 접근성이 낮았던 남부와 서부에서 특히 높았습니다.
\end{itemize}
\textbf{시사점:} 캡티브 금융 회사는 은행 신용 접근성이 낮은 지역에서 특히 강력한 역할을 수행했습니다.
\end{tcolorbox}

은행과 달리 캡티브 금융 회사는 예금을 받지 않습니다. 대신 캡티브 금융 회사는 '자산유동화기업어음(Asset-Backed Commercial Paper, ABCP)'과 '자산유동화증권(Asset-Backed Securities, ABS)'을 발행하여 자금을 조달했습니다.

\paragraph{자산유동화기업어음 (ABCP)의 작동 방식}
이러한 맥락에서 자산유동화기업어음이 어떻게 작동했는지 먼저 논의해 봅시다. 캡티브 금융 회사는 신탁(trust)을 설립하고 자동차 대출을 그 신탁으로 이전합니다. 그러면 신탁은 자동차 대출과 이 대출에서 발생하는 현금 흐름을 담보로 하는 기업어음(commercial paper)을 발행합니다. 신탁은 자산유동화기업어음 발행 수익금을 사용하여 캡티브 금융 회사에 대출금을 지급합니다.

\begin{tcolorbox}[title=\textbf{ABCP의 회전식(Revolving) 구조}]
\textbf{핵심 내용:}
\begin{itemize}
    \item 일반적으로 이러한 신탁은 '회전식(revolving)'입니다. 즉, 오래된 자동차 대출이 상환되면 캡티브 금융 회사는 새로운 자동차 대출을 신탁에 넣습니다.
    \item 이를 통해 캡티브 금융 회사는 지속적으로 단기 자금 조달에 접근할 수 있으며, 이는 다시 새로운 대출을 발행할 수 있게 합니다.
    \item 그러나 기업어음은 만기가 최대 270일이며 지속적으로 롤오버(roll over)되어야 합니다.
\end{itemize}
\textbf{시사점:} 캡티브 금융 회사는 은행과 마찬가지로 단기적으로 차입하고 장기적으로 대출합니다. 즉, '만기 변환(maturity transformation)'에 참여합니다. 이는 뱅크런과 유사한 위험에 노출될 수 있음을 의미합니다.
\end{tcolorbox}

\paragraph{자산유동화증권 (ABS)의 작동 방식}
자산유동화증권은 약간 다릅니다. 여기서는 캡티브 금융 회사가 신탁을 설정하고 자동차 대출을 그 신탁으로 이전합니다. 신탁은 더 장기적인 채권(notes)을 발행하여 자금을 조달합니다.

\begin{tcolorbox}[title=\textbf{ABS의 특징 (예: Ford Motor Credit Owner Trust 2018-B)}]
\textbf{핵심 내용:}
\begin{itemize}
    \item 포드 모터 크레딧 오너 트러스트 2018-B의 예시에서, 거의 10억 달러에 달하는 채권을 발행합니다.
    \item 이 채권은 포드 모터 크레딧이 구매한 자동차, 경트럭, 유틸리티 차량 대출 풀을 담보로 합니다.
    \item 최대 만기는 7년으로, 훨씬 더 긴 자금 조달 기간입니다. 이 거래에서는 만기 변환이 많지 않습니다.
    \item 만기가 길수록 채권의 이자율이 높아집니다.
    \item ABCP와 달리, 이것은 일회성 거래이며 대출 풀이 보충되지 않고 대출 상환금이 채권 상환에 사용됩니다.
\end{itemize}
\textbf{시사점:} 2008년 금융 위기 이전에는 ABCP가 ABS보다 저렴하고 회전식이었기 때문에 캡티브 금융 회사의 주요 자금 조달원이었습니다.
\end{tcolorbox}

\subsubsection*{2008년 금융 위기 중 ABCP 런}
그러나 2006년과 2007년에 서브프라임 위기가 모기지 시장을 강타하기 시작하자, 투자자들은 자산유동화기업어음을 담보하는 대출 풀의 품질에 대해 의심을 품기 시작했습니다. 제기된 질문 중 하나는 해당 대출 풀에 얼마나 많은 서브프라임 자동차 대출이 포함되어 있는가였습니다. 투자자들이 제너럴 모터스 어셉턴스 코퍼레이션(General Motors Acceptance Corporation, GMAC)의 풀이 낮은 품질이라는 것을 발견하자, 이러한 우려는 자동차 대출 시장으로 확산되었습니다.

투자자들은 자동차 대출을 담보로 하는 자산유동화기업어음을 롤오버하려는 의지가 점점 줄어들었고, 이는 캡티브 금융 회사에 가용 자금을 감소시켰습니다. GMAC의 대응은 극적이었습니다. 2008년 10월, GMAC는 신용 기준을 상당히 강화하고 신용 점수가 700점 이상인 소비자에게만 대출할 것이라고 발표했습니다. 당시 중간 신용 점수는 720점이었습니다. 이는 잠재적인 자동차 구매자 중 상당수가 신용을 거부당할 것이라는 의미였습니다. 다른 캡티브 금융 회사들도 이미 신용 기준을 강화했습니다. 그럼에도 불구하고 투자자들은 자산유동화기업어음 롤오버를 거부했고, 이 자금 조달은 완전히 고갈되었습니다.

신탁의 붕괴는 캡티브 금융 회사들이 자산유동화기업어음을 상환해야 한다는 것을 의미했습니다. GMAC와 크라이슬러 파이낸셜 인코퍼레이티드(Chrysler Financial Incorporated)는 의무를 이행할 수 없었고 붕괴되었습니다. 투자자들이 자산유동화기업어음 롤오버를 거부한 것은 본질적으로 뱅크런과 동일합니다. 모든 투자자들이 동시에 단기 자금 조달을 철회하여 신탁의 붕괴와 청산으로 이어지는 것입니다. 자산유동화기업어음 신탁에 대한 이러한 런과 그로 인한 신용 경색이 자동차 대출에 어떤 의미를 가졌는지 살펴보겠습니다.

\begin{tcolorbox}[title=\textbf{자동차 대출 및 판매에 미친 영향}]
\textbf{차트 분석 (텍스트 기반 설명):}
\begin{itemize}
    \item \textbf{자동차 대출 성장률:} 2006년 초 모기지 시장에서 서브프라임 위기가 전개되기 시작하면서 자동차 대출의 전년 대비 성장률은 마이너스로 전환되었습니다.
    \item \textbf{ABCP 런 이후:} 2008년 중반 투자자들이 ABCP 롤오버를 중단하자, 대출 발생은 전년 대비 30\% 이상 급감했습니다.
    \item \textbf{자동차 판매:} 금융 회사에 대한 자금 조달이 고갈되기 전에는 미국 가구가 연간 약 1,600만 대의 신차를 구매했습니다. 금융 위기가 닥치자 2009년 1분기에는 연간 790만 대로 소매 자동차 판매가 붕괴되었습니다.
\end{itemize}
\textbf{수요 감소 vs. 신용 공급 감소:}
\begin{itemize}
    \item 타이밍으로 볼 때 자동차 판매 감소는 신용 공급 감소에 의해 주도되었음을 시사합니다.
    \item 하지만 수요 감소 때문이 아니었을까요? 금융 위기 동안 많은 노동자들이 일자리를 잃었고, 암울한 경제 전망은 소비자들이 신차 구매를 주저하게 만들었을 수 있습니다.
    \item 만약 붕괴가 수요에 의해 주도되었다면, 캡티브 금융 회사들이 대출을 할 수 없을 때 다른 대출 기관들도 대출을 하지 않았을 것입니다.
    \item 그러나 2008년에는 더 많은 자동차 구매자들이 현금으로 지불하거나 은행 또는 신용 조합에서 신용을 얻었습니다. 이는 ABCP 런이 자동차 판매에 상당한 영향을 미 미쳤음을 시사합니다.
\end{itemize}
\textbf{결론:} 캡티브 금융 회사의 사례는 뱅크런이 은행 부문 밖에서도 발생할 수 있음을 보여줍니다. 실제로 만기 변환에 참여하는 모든 금융 기관은 런의 대상이 될 수 있습니다.
\end{tcolorbox}

\subsubsection*{비은행 금융 기관의 규제 공백}
그러나 2008년 금융 위기 이전에 규제 당국은 비은행 금융 기관의 만기 변환 및 신용 제공을 대체로 간과했습니다. 이러한 비은행 금융 기관은 단기 자금 시장의 런을 막는 데 도움이 될 최종 대부자로서의 연방준비제도에 접근할 수 없었습니다. 이러한 런은 신용 공급에 미치는 영향을 통해 뒤따른 경기 침체의 깊이에 기여했습니다.

\subsubsection*{2008년 금융 위기 중 자동차 대출 사례의 교훈}
이 강의에서 무엇을 배웠습니까?

\begin{tcolorbox}[title=\textbf{자동차 대출 사례의 주요 교훈}]
\begin{enumerate}
    \item \textbf{비은행 금융 기관의 만기 변환:} 캡티브 금융 회사를 예로 들어, 비은행 금융 기관도 만기 변환에 참여한다는 것을 보여주었습니다.
    \item \textbf{비은행 금융 기관의 런:} 비은행 금융 기관은 2008년 금융 위기 동안 자금 시장에서 런을 경험했습니다.
    \item \textbf{신용 공급 감소의 영향:} 런은 자동차 신용을 감소시켰고, 따라서 자동차 판매를 감소시켰습니다.
\end{enumerate}
\end{tcolorbox}

\newpage

\section{Lesson 4-2: 금융 위기에 대한 정책 대응}

\subsection{Lesson 4-2.1: 최종 대부자로서의 연준}
최종 대부자(Lender of Last Resort) 기능에 대한 수업에 오신 것을 환영합니다. 1907년 공황, 대공황, 2008년 글로벌 금융 위기 등 은행 공황과 금융 위기는 금융 부문과 궁극적으로 실물 경제에 파괴적인 영향을 미칠 수 있습니다. 이 세션에서는 중앙은행이 금융 위기 동안 단기 현금 대출을 제공함으로써 금융 시장을 안정화하고 최악의 시나리오를 방지하는 데 어떻게 도움이 될 수 있는지 살펴볼 것입니다. 이러한 개입이 원칙적으로 어떻게 설계되어야 하는지, 그리고 실제로 어떻게 작동하는지 보게 될 것입니다. 또한 이러한 개입의 잠재적 비용도 이해할 것입니다. 우리는 2008년 금융 위기와 2020년 글로벌 팬데믹의 맥락에서 이러한 문제들을 검토할 것입니다.

\subsubsection*{월터 배젓의 원칙}
정부 지원을 받는 중앙은행이 금융 위기에 개입하여 위기를 막아야 한다는 생각은 19세기 영국 언론인 월터 배젓(Walter Bagehot)에게서 시작되었습니다.

1873년, 배젓은 금융 위기 시 중앙은행은 '건전한 예금 기관에 자유롭게 대출해야 하지만, 건전한 담보에 대해서만, 그리고 진정으로 필요한 차입자가 아닌 이들을 단념시킬 만큼 충분히 높은 이자율로 대출해야 한다'고 유명하게 말했습니다. 이 진술은 중앙은행이 금융 위기에 어떻게 대응해야 하는지에 대한 시금석이 되었습니다. 이것이 무엇을 의미하는지 빠르게 풀어봅시다.

\paragraph{건전한 예금 기관에 자유롭게 대출}
"중앙은행은 건전한 예금 기관에 자유롭게 대출해야 한다"는 부분부터 시작해 봅시다. 때때로 건전한 은행들도 은행 공황에 휩쓸립니다. 모든 은행은 단기적으로 차입하고 장기적으로 대출하기 때문에, 이러한 건전한 은행 중 일부는 현금이 바닥날 수 있습니다. 중앙은행이 지원해야 하는 대상은 바로 이러한 '건전한(solvent)' 은행들입니다.

\paragraph{건전한 담보와 충분히 높은 이자율}
두 번째 부분은 어떻습니까? "건전한 담보에 대해서만, 그리고 진정으로 필요한 차입자가 아닌 이들을 단념시킬 만큼 충분히 높은 이자율로 대출해야 한다." 이것은 무엇을 달성할까요?

\begin{tcolorbox}[title=\textbf{월터 배젓 원칙의 두 가지 목표}]
\begin{enumerate}
    \item \textbf{납세자 손실 방지:} 첫째, 중앙은행은 개입할 때 손실을 입어서는 안 됩니다. 이러한 프로그램은 납세자에게 손실을 초래해서는 안 됩니다.
    \item \textbf{일시적 지원:} 둘째, 차입자들은 다른 선택지가 없을 때만 중앙은행의 도움을 받아야 합니다. 이는 일시적인 지원이어야 하며 장기적인 해결책이 되어서는 안 됩니다. 차입자들은 가능한 한 빨리 중앙은행 대출을 상환하기를 원해야 합니다.
\end{enumerate}
\textbf{결론:} 배젓은 정부가 개입해야 하지만 신중하게 해야 한다고 생각했습니다. 궁극적인 목표는 납세자에게 손실을 주지 않고 은행 공황을 막는 것입니다.
\end{tcolorbox}

이것이 이론입니다. 이제 이것이 실제로 어떻게 작동하는지 살펴보겠습니다. 최근 두 가지 위기, 즉 2008년 금융 위기와 2020년 글로벌 팬데믹에서 중앙은행의 개입을 고려해 봅시다.

\subsubsection*{2008년 금융 위기 대응}
2008년 사건부터 시작하겠습니다. 이 위기는 2006년에 시작된 미국 주택 시장의 위축으로 촉발되었습니다. 주택 가격 하락은 모기지 연체율 상승을 동반했습니다. 손실이 증가하면서 신용 공급, 특히 단기 자금 시장에서 급격한 감소가 있었습니다. 이는 유명한 투자 은행인 리먼 브라더스(Lehman Brothers)의 붕괴로 절정에 달했으며, 이는 글로벌 금융 시장에 충격을 주었습니다. 2008년 가을까지 미국은 깊은 경기 침체 직전에 있었습니다.

\paragraph{연방준비제도의 대응}
연방준비제도는 어떻게 대응했을까요? 첫 번째로 한 일은 2007년 중반에 금리를 5.25\%에서 사실상 0\%로 인하한 것이었습니다. 또한 연방준비제도는 금융 기관과 신용 시장을 안정화하기 위해 다양한 비상 유동성 프로그램(emergency liquidity programs)을 설정했습니다. 연준은 리먼 브라더스 파산 이후의 공황을 막고 싶었습니다. 이것이 연준이 최종 대부자로서 행동한 것입니다.

전통적으로 연방준비제도는 '할인 창구(discount window)'를 통해 은행에만 비상 현금 대출을 제공했을 것입니다. 그러나 2008년에는 이것이 작동하지 않았습니다. 은행만이 유일한 플레이어가 아니었습니다. 비은행 금융 기관(non-banks)이 많은 자금 시장에서 중요한 역할을 했고, 이들 비은행 기관 중 다수가 공황에 휩쓸렸습니다. 예를 들어, 리먼 브라더스는 투자 은행이었습니다. 할인 창구에 접근할 수 없었습니다.

자금 시장을 안정화하기 위해 연방준비제도는 세 가지 차원에서 할인 창구의 범위를 확장함으로써 단호하게 행동했습니다.

\begin{tcolorbox}[title=\textbf{2008년 연준의 할인 창구 확장 프로그램}]
\begin{enumerate}
    \item \textbf{TAF (Term Auction Facility):} 은행의 장기 자금 조달이 고갈되자 연방준비제도는 TAF를 설립했습니다. 이는 은행에 장기 유동성을 경매 방식으로 제공하여 자금 조달 압박을 완화했습니다.
    \item \textbf{TALF (Term Asset-backed-securities Loan Facility):} 리먼 브라더스 파산 이후 증권화 시장이 붕괴되자 연방준비제도는 이 시장을 해동하기 위해 TALF를 설립했습니다. TALF 참여자들은 연준으로부터 직접 장기 대출을 받아 증권화된 대출을 구매할 수 있었습니다.
    \item \textbf{PDCF (Primary Dealer Credit Facility) \& TSLF (Term Securities Lending Facility):} 자금 시장의 핵심인 많은 비은행 기관들이 절실히 현금을 필요로 했습니다. 가장 중요하고 평판 좋은 투자 은행들은 '프라이머리 딜러(primary dealers)'라고도 불립니다. 프라이머리 딜러는 연방준비제도의 공개 시장 운영의 거래 상대방 역할을 합니다. 이들은 통화 정책에 매우 중요합니다. 또한 잘 기능하는 국채 시장에도 필수적입니다. PDCF와 TSLF는 이들의 필요를 충족시키기 위해 설립되었습니다.
    \item \textbf{AMLF (Asset-Backed Commercial Paper Money Market Mutual Fund Liquidity Facility):} 머니 마켓 펀드(money market funds)가 런을 경험하고 있었기 때문에 연준은 AMLF를 설립했습니다. 이는 머니 마켓 펀드에 유동성을 제공하여 투자자들의 대규모 인출을 막았습니다.
    \item \textbf{CPFF (Commercial Paper Funding Facility):} 마지막으로 CPFF는 미국 상업어음 발행자들에게 지원을 제공했습니다. 이것은 2008년의 유일한 직접 구매 프로그램의 예입니다. 이는 연방준비제도가 상업어음을 직접 구매했다는 것을 의미합니다. 이것은 주로 금융 기관이 발행한 것이었지만, 연준은 할리 데이비슨(Harley Davidson)이나 맥도날드(McDonald's)와 같은 비금융 기관이 발행한 어음도 구매했습니다. 역사적 기준으로 볼 때, 이것은 매우 이례적인 조치였습니다.
\end{enumerate}
\end{tcolorbox}

\paragraph{2008년 개입의 성공 여부와 비용}
그것은 크고 복잡한 개입이었습니다. 성공적이었을까요? 월터 배젓은 만족했을까요?

\begin{tcolorbox}[title=\textbf{2008년 연준 개입의 평가}]
\textbf{긍정적 측면:}
\begin{itemize}
    \item \textbf{금융 공황 중단:} 무엇보다도 연방준비제도는 금융 공황을 막았습니다.
    \item \textbf{손실 없음:} 각 프로그램에서 비상 대출은 고품질 담보로 뒷받침되었고 시장 금리 이상의 이자율이 부과되었습니다. 대출은 신속하게 상환되었고, 프로그램은 손실 없이 종료되었습니다.
\end{itemize}
\textbf{잠재적 비용 및 우려:}
\begin{itemize}
    \item \textbf{정부 안전망의 확장:} 다른 한편으로, 이것은 정부 안전망의 대규모적이고 전례 없는 확장(unprecedented expansion)이었습니다. 지원을 받은 금융 기관 중 다수는 예금 기관이 아니었습니다. 그들은 연방준비제도나 다른 은행 규제 기관의 감독이나 규제를 받지 않았습니다.
    \item \textbf{도덕적 해이(Moral Hazard) 우려:} 만약 시장 참여자들이 상황이 나빠질 때 연방준비제도가 계속해서 지원을 제공할 것이라고 기대한다면, 이것이 위험 감수 행동에 어떤 영향을 미칠까요? 이 질문에 대한 명확한 답은 없습니다.
\end{itemize}
\end{tcolorbox}

\subsubsection*{2020년 글로벌 팬데믹 대응}
2020년으로 빠르게 넘어가 봅시다. 글로벌 팬데믹은 2020년 1분기에 미국에 도착했습니다. 2020년 3월까지 국채 시장에 심각한 변동성이 있었고, 이는 단기 자금 시장에 혼란을 야기할 가능성이 있었습니다. 이는 전 세계적으로 경고음을 울렸습니다. 2008년 금융 위기가 얼마나 심각해질 수 있는지를 보여주었기 때문에, 정책 입안자들은 국채 시장에 대한 첫 번째 지원을 제공하고 더 광범위하게 지원하기 위해 서둘렀습니다. 이전과 마찬가지로, 여기에는 최종 대부자로서 행동하는 연방준비제도가 포함되었습니다.

충격은 매우 달랐지만, 중앙은행의 대응은 유사했습니다. 2008년 금융 위기 동안 사용되었던 여러 유동성 시설이 다시 가동되었습니다. 여기에는 TALF(Term-Asset-Backed Securities Loan Facility), CPFF(Commercial Paper Funding Facility), MMLF(Money Market Mutual Fund Liquidity Facility)가 포함되었습니다.

\begin{tcolorbox}[title=\textbf{2020년 연준의 주요 유동성 프로그램}]
\begin{enumerate}
    \item \textbf{PPPLF (Paycheck Protection Program Liquidity Facility):} 아마도 가장 큰 프로그램은 PPPLF였을 것입니다. 급여 보호 프로그램(Paycheck Protection Program)은 중소기업에 6천억 달러의 상환 가능한 대출을 제공하는 정부 프로그램이었습니다. PPPLF에 따라 대출 기관은 중앙은행으로부터 이러한 대출에 대한 자금을 받을 수 있었습니다.
    \item \textbf{PSMCCF (Primary and Secondary Market Corporate Credit Facilities):} 직접 구매 프로그램은 이번에는 훨씬 더 큰 역할을 했습니다. PSMCCF는 회사채 시장을 통해 대기업에 직접적인 지원을 제공했습니다. 연준은 대규모로 회사채를 직접 구매했습니다.
    \item \textbf{Main Street Lending Program:} 2020년 4월까지 연방준비제도는 중소기업에 대한 직접 대출에도 관여했습니다. 메인 스트리트 대출 프로그램(Main Street Lending Program)은 연준이 은행과 협력하여 중소기업 대출을 시작할 수 있도록 했습니다.
\end{enumerate}
\end{tcolorbox}

\paragraph{2020년 개입의 교훈}
이것은 우리에게 무엇을 남길까요? 2020년 사건에서 최종 대부자에 대해 무엇을 더 배울 수 있을까요?

\begin{tcolorbox}[title=\textbf{2020년 연준 개입의 평가 및 교훈}]
\textbf{긍정적 측면:}
\begin{itemize}
    \item \textbf{경제 붕괴 방지:} 이러한 개입은 경제의 완전한 붕괴를 피하는 데 성공한 것으로 보입니다.
    \item \textbf{재정-통화 정책 협력:} 위기 개입은 재정 정책과 통화 정책이 협력할 때 가장 잘 작동하는 것으로 보입니다. 아마도 가장 성공적인 정부 개입은 급여 보호 프로그램이었을 것입니다. 연방준비제도가 이러한 대출에 6천억 달러의 자금을 제공한 것이 위기 시기에 얼마나 중요했는지 쉽게 알 수 있습니다. 이 시점에서 은행들은 추가 대출을 하기보다는 현금을 비축했을 것입니다.
\end{itemize}
\textbf{월터 배젓의 관점과 우려:}
\begin{itemize}
    \item 월터 배젓은 이 모든 것에 대해 어떻게 생각했을까요? 그는 중앙은행이 예금 기관만 지원해야 한다고 생각했습니다. 그것은 과거의 일이 된 것 같습니다.
    \item 그는 또한 납세자에게 잠재적인 손실이 발생할까 봐 우려했습니다. 불행히도 2008년에서 2020년으로 가면서 연방준비제도가 중소기업 및 대기업에 대한 직접 대출에 훨씬 더 많이 관여하는 것을 보았습니다. 이는 물론 기업 부문에서 채무 불이행이 발생할 경우 연방준비제도가 직접적인 손실에 노출된다는 것을 의미합니다. 이는 분명히 배젓을 불안하게 만들 것입니다.
\end{itemize}
\end{tcolorbox}

\subsubsection*{최종 대부자 역할의 진화}
이 세션에서 무엇을 배웠습니까?

\begin{tcolorbox}[title=\textbf{최종 대부자 역할의 주요 교훈}]
\begin{enumerate}
    \item \textbf{위기 중단 능력:} 최종 대부자는 비상 대출 프로그램을 통해 금융 위기를 막을 수 있습니다.
    \item \textbf{역사적 원칙:} 역사적으로 이러한 자금은 은행으로 흘러갔고, 대출은 고품질 담보와 시장 금리 이상의 이자율을 요구했습니다.
    \item \textbf{최근 개입 사례:} 우리는 연방준비제도가 2008년과 2020년에 시장을 진정시키기 위해 어떻게 개입했는지 보았습니다.
    \item \textbf{프로그램 범위 확장:} 마지막으로, 이러한 비상 대출 프로그램의 범위가 최근 몇 년 동안 비은행 기관뿐만 아니라 직접 대출 프로그램까지 포함하도록 어떻게 확장되었는지 논의했습니다.
\end{enumerate}
\end{tcolorbox}

\newpage

\subsection{Lesson 4-2.2: 금융 위기에 대한 재정적 대응}
이 강의는 금융 위기에 대한 재정적 대응에 대해 다룹니다. (이 파트 1/2 텍스트에는 해당 강의의 내용이 포함되어 있지 않습니다. 전체 모듈 4의 파트 2/2에서 다루어질 예정입니다.)

\newpage

\newpage

\section*{개요: 재정 위기에 대한 재정 대응}

이 문서는 금융 위기와 그에 따른 경기 침체에 정부가 어떻게 재정 정책(Fiscal Policy)으로 대응하는지 설명합니다. 통화 정책(Monetary Policy)만으로는 경제를 완전히 회복시키기 어렵기 때문에, 정부는 다양한 개입을 통해 경제 회복을 지원합니다. 우리는 세 가지 주요 재정 대응 방안을 살펴보고, 각 방안의 효과를 2008년 금융 위기와 2020년 팬데믹 사례를 통해 분석합니다.

\begin{itemize}
    \item \textbf{주요 목표:} 금융 위기로 인한 경제적 피해를 줄이고 경기 침체에서 벗어나기 위한 재정 정책의 역할 이해.
    \item \textbf{핵심 개념:} 정부의 소비자 직접 이전(세금 감면, 경기 부양 수표), 기업 지원, 직접 지출(인프라 프로젝트).
    \item \textbf{평가 포인트:} 각 재정 대응 방안의 효과성, 케인즈주의적 관점과 반대 관점, 통화 정책의 보조적 역할.
\end{itemize}

\newpage
\section*{용어 정리}

다음은 본 문서에서 사용되는 주요 경제 용어와 그에 대한 쉬운 설명입니다.

\begin{adjustbox}{width=\textwidth,center}
\begin{tabular}{lllc}
    \toprule
    \textbf{용어} & \textbf{쉬운 설명} & \textbf{원어} & \textbf{비고} \\
    \midrule
    재정 정책 & 정부가 세금과 지출을 조절하여 경제를 안정시키려는 정책 & Fiscal Policy & 통화 정책과 대비 \\
    통화 정책 & 중앙은행이 이자율과 통화량을 조절하여 경제를 안정시키려는 정책 & Monetary Policy & 재정 정책과 대비 \\
    금융 위기 & 금융 시스템이 불안정해져 경제 전반에 큰 충격을 주는 상황 & Financial Crisis & 2008년, 2020년 사례 \\
    경기 침체 & 경제 활동이 전반적으로 위축되고 생산과 고용이 감소하는 시기 & Recession & GDP 감소, 실업률 증가 \\
    GDP & 한 나라에서 일정 기간 생산된 모든 최종 재화와 서비스의 시장 가치 합계 & Gross Domestic Product & 경제 규모를 나타내는 지표 \\
    소비 (C) & 가계가 재화와 서비스를 구매하는 지출 & Consumption & GDP의 가장 큰 구성 요소 \\
    투자 (I) & 기업이 생산 설비 등에 지출하거나 가계가 주택을 구매하는 지출 & Investment & 미래 생산 능력 증대 \\
    정부 지출 (G) & 정부가 재화와 서비스를 구매하는 지출 & Government Spending & 인프라 건설, 공무원 급여 등 \\
    순수출 (NX) & 수출에서 수입을 뺀 값 & Net Exports & (수출 - 수입) \\
    경기 부양 수표 & 정부가 경제 활성화를 위해 가계에 직접 지급하는 현금 & Stimulus Checks & 소비 진작 목적 \\
    실업 급여 & 실직자에게 정부가 지급하는 소득 보전금 & Unemployment Benefits & 소비 감소 방지 \\
    임금 보조금 & 정부가 기업의 임금 일부를 지원하여 해고를 막는 정책 & Wage Bill Subsidies & 고용 유지 목적 \\
    급여 보호 프로그램 & 팬데믹 시기 미국에서 기업의 고용 유지를 위해 제공된 대출 프로그램 & Paycheck Protection Program (PPP) & 대출 상환 면제 조건 \\
    정부 지출 승수 & 정부 지출 1달러가 GDP를 얼마나 증가시키는지 나타내는 값 & Government Spending Multiplier & 케인즈주의 이론의 핵심 \\
    \bottomrule
\end{tabular}
\end{adjustbox}

\newpage
\section*{핵심 개념 및 원리: 재정 위기 대응 방안}

금융 위기와 경기 침체 시, 경제를 회복시키기 위해 정부는 재정 정책을 활용합니다. 통화 정책만으로는 한계가 있기 때문에, 정부는 다양한 개입을 통해 경제 회복을 지원합니다.

\subsection*{1. GDP의 구성 요소와 소비의 중요성}

\begin{tcolorbox}[mybox={GDP의 구성 요소}]
    \textbf{한 줄 핵심 요약:} GDP는 소비, 투자, 정부 지출, 순수출의 합으로, 이 중 소비가 가장 큰 비중을 차지합니다.
\end{tcolorbox}

경제의 총 생산량, 즉 국내총생산(GDP)은 항상 다음 네 가지 요소로 구성됩니다:
\begin{itemize}
    \item \textbf{민간 소비 (C, Private Consumption):} 가계가 재화와 서비스를 구매하는 지출입니다.
    \item \textbf{기업 투자 (I, Business Investment):} 기업이 공장, 기계 등 생산 설비에 지출하거나 가계가 주택을 구매하는 지출입니다.
    \item \textbf{정부 지출 (G, Government Spending):} 정부가 재화와 서비스를 구매하는 지출입니다.
    \item \textbf{순수출 (NX, Net Exports):} 수출에서 수입을 뺀 값입니다. (수출 - 수입)
\end{itemize}

\begin{examplebox}{예시: 2019년 미국 GDP 구성}
    2019년 미국의 경우, GDP에서 소비가 약 68\%를 차지했습니다. 기업 투자는 17\%, 정부 지출은 19\%, 그리고 순수출은 -4\%였습니다. 순수출이 마이너스라는 것은 미국이 수출보다 수입을 더 많이 했다는 의미입니다. 이처럼 소비는 GDP에서 가장 큰 비중을 차지하는 요소입니다.
\end{examplebox}

\begin{tcolorbox}[mybox={경기 침체 시 소비의 변화}]
    \textbf{한 줄 핵심 요약:} 경기 침체 시 소비는 급격히 감소하며, 이는 GDP에 막대한 영향을 미칩니다.
\end{tcolorbox}

경기가 좋을 때는 소비가 꾸준히 증가하지만, 금융 위기나 팬데믹과 같은 어려운 시기에는 소비가 급격히 감소합니다. 예를 들어, 2008년 금융 위기와 그 여파, 그리고 2020년 글로벌 팬데믹 시기에는 소비가 크게 줄었습니다. 이 두 경우 모두 실업률이 급증했고, 이는 가계의 소득 감소로 이어져 소비를 더욱 위축시켰습니다. 소비가 GDP의 3분의 2를 차지하기 때문에, 소비의 급락은 경제 전반에 엄청난 충격을 줍니다. 따라서 재정 위기에 대한 정부의 대응은 주로 소비를 지탱하는 데 초점을 맞춥니다.

\subsection*{2. 재정 대응 방안 1: 소비자에게 직접 이전 (Government Transfers to Consumers)}

정부가 소비를 지탱하기 위해 가장 먼저 고려하는 방법 중 하나는 가계에 직접 돈을 이전하는 것입니다.

\subsubsection*{2.1. 세금 감면 vs. 경기 부양 수표}

\begin{tcolorbox}[mybox={세금 감면의 한계}]
    \textbf{한 줄 핵심 요약:} 실업자에게는 소득이 없으므로, 소득세 감면은 경기 침체 시 소비 진작에 효과적이지 않습니다.
\end{tcolorbox}

새롭게 실업자가 된 사람들은 소득이 없기 때문에, 소득세를 감면해주는 정책은 이들에게 아무런 도움이 되지 않습니다. 이들은 소비를 가장 많이 줄이는 계층이므로, 이들을 직접 지원하는 다른 방법이 필요합니다.

\begin{tcolorbox}[mybox={경기 부양 수표의 도입}]
    \textbf{한 줄 핵심 요약:} 정부는 모든 가구에 경기 부양 수표를 보내 소비를 직접적으로 늘리려 합니다.
\end{tcolorbox}

세금 감면의 한계를 보완하기 위해 정부는 모든 가구에 경기 부양 수표(Stimulus Checks)를 보냅니다. 이는 가계에 직접 현금을 지급하여 소비를 유도하는 방식입니다.

\begin{examplebox}{예시: 미국 경기 부양 수표 지급 사례}
    \begin{itemize}
        \item \textbf{2008년:} 개인 납세자는 최대 600달러, 부부는 최대 1,200달러를 받았습니다.
        \item \textbf{2020년:} 개인 납세자는 최대 1,200달러, 부부는 최대 2,400달러를 받았습니다.
    \end{itemize}
\end{examplebox}

\begin{tcolorbox}[mybox={경기 부양 수표의 효과성}]
    \textbf{한 줄 핵심 요약:} 경기 부양 수표는 소비 진작에 부분적으로 효과적이지만, 모든 금액이 소비로 이어지지는 않습니다.
\end{tcolorbox}

경기 부양 수표가 과연 효과적일까요? 연구 결과에 따르면, 2020년 지급된 경기 부양 수표의 경우 약 42\%만이 소비에 사용되었고, 27\%는 저축되었으며, 31\%는 부채 상환에 사용되었습니다.

\begin{cautionbox}{왜 모든 돈이 소비되지 않았을까요?}
    \begin{itemize}
        \item \textbf{소득 수준에 따른 차이:} 저소득층은 수표의 대부분을 소비했지만, 고소득층은 주로 저축했습니다. 이는 고소득층이 이미 충분한 소비 여력이 있거나, 미래의 불확실성에 대비해 저축을 선호하기 때문입니다.
        \item \textbf{소비 기회의 제한:} 2020년에는 미국 경제가 봉쇄(Lockdown)되어 소비할 기회가 제한적이었습니다. 상점들이 문을 닫고 외식이나 여행이 어려워지면서, 돈을 쓰고 싶어도 쓸 곳이 마땅치 않았던 것입니다.
    \end{itemize}
\end{cautionbox}

\subsubsection*{2.2. 목표 지향적 직접 지급: 실업 급여 증액 및 고용 유지 프로그램}

\begin{tcolorbox}[mybox={실업 급여의 효과성}]
    \textbf{한 줄 핵심 요약:} 실업 급여 증액은 실직 가구의 소비 감소를 직접적으로 막아주어 매우 효과적인 개입입니다.
\end{tcolorbox}

경기 부양 수표보다 더 목표 지향적인(Targeted) 직접 지급 방식은 실업 급여(Unemployment Benefits)를 늘리는 것입니다.
\begin{itemize}
    \item \textbf{2009년:} 실업 급여 지급 기간이 연장되었습니다.
    \item \textbf{코로나19 팬데믹 시기:} 실업 급여가 증액되고 지급 기간이 다시 연장되었습니다.
\end{itemize}
이러한 개입은 최근 일자리를 잃은 사람들에게만 지급되므로, 소비를 가장 많이 줄이는 가구에 직접적인 도움을 줍니다. 이는 경기 부양 수표보다 효율적일 수 있습니다.

\begin{tcolorbox}[mybox={고용 유지 프로그램의 장점}]
    \textbf{한 줄 핵심 요약:} 기업이 직원을 해고하지 않도록 임금의 일부를 지원하는 프로그램은 고용 관계를 유지하고, 직원의 기술 손실을 막으며, 경제 회복 시 재고용 비용을 줄여줍니다.
\end{tcolorbox}

또 다른 접근 방식은 기업에 돈을 지급하여 직원을 해고하지 않도록 하는 것입니다. 정부가 기업의 임금 청구서(Wage Bill)의 일부를 부담하는 방식입니다.

\begin{examplebox}{예시: 독일과 미국의 고용 유지 프로그램}
    \begin{itemize}
        \item \textbf{독일 (2008년 금융 위기):} 이 접근 방식은 2008년 금융 위기 당시 독일에서 처음 시도되었습니다.
        \item \textbf{미국 (코로나19 팬데믹):} 미국의 급여 보호 프로그램(Paycheck Protection Program, PPP)은 이와 유사한 목적을 가졌습니다. 이 프로그램은 중소기업에 대출을 제공했는데, 만약 기업이 직원을 해고하지 않으면 이 대출을 탕감해 주었습니다.
    \end{itemize}
\end{examplebox}

이러한 프로그램의 주요 장점은 다음과 같습니다:
\begin{itemize}
    \item \textbf{고용 관계 유지:} 고용주와 직원 간의 관계를 그대로 유지합니다.
    \item \textbf{기술 손실 방지:} 직원의 기술이 감가상각(Depreciate)되지 않습니다. 즉, 일자리를 잃고 쉬는 동안 기술이 녹슬지 않습니다.
    \item \textbf{재고용 비용 절감:} 수요가 다시 증가했을 때 기업이 직원을 다시 고용할 필요가 없어 재고용에 드는 시간과 비용을 절약할 수 있습니다.
    \item \textbf{소비 유지:} 직원들이 소득을 유지하므로 소비를 줄일 필요가 없어집니다.
\end{itemize}

\begin{summarybox}{소비자 직접 이전의 결론}
    가계에 대한 직접 이전은 경기 침체 시 소비와 GDP를 지탱하는 데 도움이 됩니다. 특히, 새로 실직한 사람들을 대상으로 하거나 고용을 유지하는 프로그램이 가장 효과적입니다.
\end{summarybox}

\newpage
\subsection*{3. 재정 대응 방안 2: 기업 지원 (투자 촉진)}

GDP의 두 번째 주요 구성 요소는 투자(Investment)입니다.

\begin{tcolorbox}[mybox={경기 침체 시 투자의 변화}]
    \textbf{한 줄 핵심 요약:} 경기 침체 시 투자는 급격히 감소하며, 정부의 개입 옵션은 제한적입니다.
\end{tcolorbox}

투자는 2008년 금융 위기 때 급락했고, 2020년에도 다시 감소했습니다. 이러한 투자 감소에 대응하기 위해 정부는 다음과 같은 제한적인 옵션을 가지고 있습니다.

\begin{itemize}
    \item \textbf{세금 감면 (Tax Breaks):} 기업의 투자 비용을 줄여주기 위해 세금을 감면해 줄 수 있습니다.
    \item \textbf{보조금 (Subsidies):} 특정 투자 활동에 대해 직접적인 보조금을 지급할 수 있습니다.
\end{itemize}

\begin{cautionbox}{투자의 근본적인 동인}
    결국 투자는 경제의 전반적인 상태에 의해 좌우됩니다. 예를 들어, 기업 투자는 상품에 대한 수요에 의해 결정됩니다. 수요가 없으면 아무리 세금 감면이나 보조금을 줘도 기업은 투자를 늘리지 않을 것입니다. 따라서 정부의 직접적인 투자 촉진 정책은 한계가 있을 수 있습니다.
\end{cautionbox}

\subsection*{4. 재정 대응 방안 3: 정부의 직접 지출 (Direct Government Spending)}

\begin{tcolorbox}[mybox={정부 지출의 역할}]
    \textbf{한 줄 핵심 요약:} 정부는 인프라 프로젝트와 같은 직접 지출을 늘려 일자리를 창출하고 경제 활동을 촉진할 수 있습니다.
\end{tcolorbox}

정부가 경기 침체에 대응하는 마지막 유형의 재정 정책은 정부 자체의 지출을 늘리는 것입니다.

\begin{examplebox}{예시: 미국 정부의 직접 지출 프로그램}
    \begin{itemize}
        \item \textbf{2009년 미국 경기 회복 및 재투자법 (American Recovery and Reinvestment Act of 2009):} 이 법안은 2,750억 달러를 연방 계약, 보조금, 대출에 할당하여 일자리를 창출했습니다.
        \item \textbf{2021년 인프라 투자 및 일자리 법 (Infrastructure Investment and Jobs Act):} 글로벌 팬데믹 이후 통과된 이 법안은 전국적으로 많은 공공 건설 프로젝트를 포함했습니다.
    \end{itemize}
\end{examplebox}

이러한 정부 지출 법안들은 주로 도로, 교량, 공항 등 공공 인프라 건설 프로젝트를 포함합니다. 하지만 정부 지출을 늘리는 것이 경기 침체에 맞서는 효과적인 도구일까요?

\subsubsection*{4.1. 정부 지출 승수 (Government Spending Multiplier)}

\begin{tcolorbox}[mybox={케인즈주의적 관점}]
    \textbf{한 줄 핵심 요약:} 케인즈는 경기 침체 시 정부 지출이 GDP를 1달러 이상 증가시키는 승수 효과를 가진다고 주장했습니다.
\end{tcolorbox}

영국의 경제학자 존 메이너드 케인즈(John Maynard Keynes)는 경기 침체 시 소비가 감소하면 정부가 직접 상품을 구매하여 이 손실된 수요를 보충해야 한다고 주장했습니다. 케인즈는 이것이 GDP를 안정화시킬 것이라고 보았습니다. 여기서 핵심 질문은 "정부가 1달러를 지출했을 때 GDP가 1달러보다 더 많이 증가하는가?"입니다. 이것을 \textbf{정부 지출 승수(Government Spending Multiplier)}라고 부릅니다.

케인즈주의적 관점에서는 경기 침체 시 정부 지출 승수가 1보다 크다고 봅니다. 그 이유는 정부가 손실된 수요를 보충하면 일자리가 창출되고, 이는 다시 소비를 증가시키기 때문입니다. 즉, 정부가 1달러를 지출하면 그 돈이 누군가의 소득이 되고, 그 소득의 일부가 다시 소비되어 또 다른 사람의 소득이 되는 연쇄 효과가 발생한다는 것입니다.

\begin{tcolorbox}[mybox={케인즈주의 반대 관점}]
    \textbf{한 줄 핵심 요약:} 반대론자들은 정부 지출이 민간 수요를 단순히 대체할 뿐이며, 비효율적인 자원 재배치를 초래하여 승수 효과가 1보다 작다고 주장합니다.
\end{tcolorbox}

케인즈주의에 반대하는 사람들은 정부 지출 승수가 1보다 작다고 주장합니다. 이들은 정부가 단순히 민간 소비를 대체할 수 없다고 봅니다. 정부는 민간이 원하지 않는 다른 종류의 상품에 대한 수요를 창출할 수 있다는 것입니다.

\begin{examplebox}{예시: 비효율적인 자원 재배치}
    주택 부문에서 발생한 충격에 대응하여 정부가 군사 지출을 늘린다고 가정해 봅시다. 건설 부문의 노동자들은 국방 부문에서 일자리를 찾아야 할 것입니다. 하지만 건설 노동자들이 국방 산업에 필요한 기술을 가지고 있지 않을 수 있습니다. 노동자들의 기술을 재교육하고 산업 간에 재배치하는 것은 비용이 많이 들고 비효율적일 수 있습니다. 따라서 정부 지출 승수는 1보다 작다는 주장이 나옵니다.
\end{examplebox}

\begin{summarybox}{정부 지출 승수에 대한 증거}
    대부분의 연구는 경기 침체 시 정부 지출 승수가 1보다 크다는 것을 시사하지만, 증거는 혼합되어 있습니다. 예상대로, 그 효과는 돈이 어떻게 사용되는지에 따라 달라집니다.
\end{summarybox}

\subsubsection*{4.2. 정부 지출의 실제적 어려움}

\begin{cautionbox}{정부 지출의 시간 지연 및 측정의 어려움}
    \textbf{한 줄 핵심 요약:} 경기 침체 시 정부 지출은 즉각적으로 이루어져야 하지만, 대규모 프로젝트는 계획, 승인, 실행에 시간이 오래 걸리며, 그 효과를 측정하기도 어렵습니다.
\end{cautionbox}

경기 침체 시 정부는 즉시 지출해야 하지만, 실제로는 프로젝트에 시간이 걸리는 경우가 많습니다.
\begin{itemize}
    \item \textbf{시간 지연:} 대규모 인프라 프로젝트를 생각해 보세요. 신중한 계획, 이해관계자들의 동의, 정부 구매 절차를 거쳐야 합니다. 이 모든 과정은 프로그램의 실행을 지연시킬 가능성이 큽니다.
    \item \textbf{효과 측정의 어려움:} 이러한 지출의 효과를 측정하는 것도 어렵습니다. 예를 들어, 수리된 다리의 이점을 어떻게 측정할 수 있을까요? 교통량 증가, 안전 개선, 경제 활동 촉진 등 다양한 측면을 고려해야 하지만, 이를 정량화하기는 쉽지 않습니다.
\end{itemize}

\begin{summarybox}{정부 지출의 결론}
    결론적으로, 특정 조건 하에서는 정부 지출을 늘리는 것이 경기 침체 시 경제를 지원하는 효과적인 방법이 될 수 있습니다.
\end{summarybox}

\newpage
\section*{통화 정책의 보조적 역할}

\begin{tcolorbox}[mybox={재정 정책을 지원하는 통화 정책}]
    \textbf{한 줄 핵심 요약:} 중앙은행은 이자율을 낮추고 금융 기관에 자금을 제공하여 정부의 재정 부양책을 지원하고 그 효과를 증폭시킬 수 있습니다.
\end{tcolorbox}

마지막으로, 통화 정책은 이러한 재정 대응을 지원하는 데 중요한 역할을 할 수 있습니다. 정부의 세 가지 유형의 개입(소비자 이전, 기업 지원, 직접 지출)은 모두 비용이 듭니다. 중앙은행은 이자율을 통제함으로써 정부의 경기 부양 프로그램을 두 가지 방식으로 지원할 수 있습니다.

\begin{enumerate}
    \item \textbf{정부의 차입 비용 직접 인하:} 중앙은행이 이자율을 낮추면 정부가 돈을 빌리는 비용이 직접적으로 줄어듭니다. 이는 정부가 더 많은 재정 지출을 할 수 있도록 재정적 여유를 제공합니다.
    \item \textbf{민간 부문 활동 증폭:} 이자율이 낮아지면 기업이 투자하고 가계가 소비를 위해 대출을 받는 것이 더 저렴해집니다. 이는 정부의 경기 부양책 효과를 증폭시킬 수 있습니다. 예를 들어, 정부가 인프라 프로젝트를 시작하면 기업은 낮은 이자율로 자금을 조달하여 프로젝트에 참여하고, 이는 다시 고용과 소비를 늘리는 선순환을 만듭니다.
\end{enumerate}

중앙은행은 또한 금융 기관에 자금을 제공하여 경기 부양 프로그램을 지원할 수 있습니다.

\begin{examplebox}{예시: 급여 보호 프로그램 유동성 지원 시설 (2020년)}
    2020년의 급여 보호 프로그램 유동성 지원 시설(Paycheck Protection Program Liquidity Facility)이 대표적인 예입니다. 이 프로그램 하에서 대출 기관들은 중앙은행으로부터 자금을 받아 정부의 급여 보호 프로그램(PPP)에 따라 중소기업에 제공되는 보조금을 지원했습니다. 중앙은행의 이러한 지원이 없었다면, 자체적인 제약에 직면한 대출 기관들은 정부 프로그램에 협조하기를 꺼렸을 수도 있습니다.
\end{examplebox}

\newpage
\section*{빠르게 훑어보기: 재정 위기 대응 요약}

\begin{summarybox}{재정 위기 대응의 핵심}
    \textbf{목표:} 금융 위기로 인한 최악의 시나리오를 피하고 경제를 지원하는 것.
\end{summarybox}

\begin{tcolorbox}[mybox={세 가지 주요 재정 대응 방안}]
    \begin{itemize}
        \item \textbf{수요 진작:}
        \begin{itemize}
            \item 가계에 직접 돈을 이전 (경기 부양 수표, 실업 급여 증액).
            \item 정부 지출 증가 (인프라 프로젝트).
        \end{itemize}
        \item \textbf{목표 지향적 개입의 효과성:}
        \begin{itemize}
            \item 임금 보조금 지급이나 실업 급여 증액과 같이 특정 계층이나 기업을 대상으로 하는 개입이 가장 효과적입니다. 이는 소비 감소를 직접적으로 막고 고용을 유지하는 데 기여합니다.
        \end{itemize}
        \item \textbf{통화 정책의 지원:}
        \begin{itemize}
            \item 중앙은행의 지원적인 통화 정책(이자율 인하, 금융 기관 유동성 공급)은 정부의 재정 부양 프로그램을 촉진하고 그 효과를 증폭시킬 수 있습니다.
        \end{itemize}
    \end{itemize}
\end{tcolorbox}

\newpage
\section*{FAQ: 자주 묻는 질문}

\begin{tcolorbox}[mybox={Q\&A: 재정 위기 대응}]
    \textbf{Q1: 통화 정책만으로는 경기 침체를 극복할 수 없나요?}
    \textbf{A1:} 통화 정책은 이자율 조절을 통해 경제를 자극할 수 있지만, 금융 위기나 심각한 경기 침체 시에는 가계와 기업의 신뢰가 낮아져 이자율 인하만으로는 투자와 소비를 충분히 늘리기 어렵습니다. 이때 정부의 직접적인 재정 개입이 필요합니다.

    \textbf{Q2: 경기 부양 수표는 왜 모든 사람에게 똑같이 효과적이지 않나요?}
    \textbf{A2:} 경기 부양 수표는 저소득층에게는 즉각적인 소비로 이어지는 경향이 강하지만, 고소득층은 이미 소비 여력이 있거나 미래 불확실성에 대비해 저축하거나 부채를 상환하는 데 사용하는 경향이 있습니다. 또한, 경제 봉쇄와 같은 상황에서는 소비할 기회 자체가 부족할 수 있습니다.

    \textbf{Q3: 정부 지출 승수가 1보다 크다는 것이 정확히 무엇을 의미하나요?}
    \textbf{A3:} 정부 지출 승수가 1보다 크다는 것은 정부가 1달러를 지출했을 때, 그 1달러가 직접적인 GDP 증가 외에 추가적인 경제 활동(예: 일자리 창출, 소득 증가, 추가 소비)을 유발하여 최종적으로 GDP가 1달러보다 더 많이 증가한다는 의미입니다. 이는 연쇄적인 경제 효과를 나타냅니다.

    \textbf{Q4: 정부 지출이 항상 좋은가요? 단점은 없나요?}
    \textbf{A4:} 정부 지출은 경기 침체 시 수요를 보충하고 일자리를 창출하는 긍정적인 효과가 있지만, 단점도 있습니다. 대규모 프로젝트는 계획 및 실행에 시간이 오래 걸려 적절한 시기를 놓칠 수 있고, 비효율적인 자원 배분으로 이어질 수 있으며, 정부 부채를 증가시키는 요인이 될 수 있습니다.

    \textbf{Q5: 중앙은행은 재정 정책을 어떻게 지원하나요?}
    \textbf{A5:} 중앙은행은 주로 이자율을 낮춰 정부의 차입 비용을 줄이고, 기업과 가계의 투자 및 소비 대출을 저렴하게 만들어 재정 부양책의 효과를 증폭시킵니다. 또한, 금융 기관에 유동성을 공급하여 정부 프로그램(예: PPP)이 원활하게 작동하도록 돕습니다.
\end{tcolorbox}

\end{document}
